% !TEX program = xelatex
\documentclass[11pt,a4paper,fontset=macnew]{ctexart}
\usepackage[left=20mm,right=20mm,top=19mm,bottom=19mm,headsep=5mm]{geometry}
\usepackage{fontspec}
\setmainfont{Times New Roman}
\setsansfont{Arial}
\usepackage{booktabs,longtable,tabularx,array,enumitem,needspace,fancyhdr,xcolor}
\usepackage[hidelinks,bookmarksopen=true,bookmarksnumbered=false]{hyperref}
\hypersetup{pdftitle={英语二作文：36类融合模板},pdfauthor={学习资料整合},pdfsubject={按用户提供的批改指令综合重写；应用文与大作文}}
\newcolumntype{L}[1]{>{\raggedright\arraybackslash}p{#1}}
\setlength{\parindent}{0pt}
\setlength{\parskip}{3pt}
\linespread{1.10}
\setlength{\emergencystretch}{2em}
\setlength{\tabcolsep}{4pt}
\renewcommand{\arraystretch}{1.18}
\setlist[itemize]{leftmargin=1.35em,itemsep=3pt,topsep=3pt}
\setlist[enumerate]{leftmargin=1.6em,itemsep=3pt,topsep=3pt}
\ctexset{section/format={\Large\bfseries},section/beforeskip=0pt,section/afterskip=8pt,subsection/format={\large\bfseries},subsection/beforeskip=9pt,subsection/afterskip=4pt}
\pagestyle{fancy}
\fancyhf{}
\fancyhead[L]{\small 英语二作文 \quad 多源融合模板}
\fancyhead[R]{\small\nouppercase{\leftmark}}
\fancyfoot[C]{\small\thepage}
\setlength{\headheight}{15pt}
\renewcommand{\headrulewidth}{0.3pt}
\newenvironment{writing}{\begingroup\fontsize{10.5}{13.5}\selectfont\setlength{\parskip}{3pt}}{\par\endgroup}
\newenvironment{example}{\begingroup\fontsize{11}{14.3}\selectfont\setlength{\parskip}{4pt}}{\par\endgroup}
\begin{document}
\thispagestyle{empty}
\vspace*{25mm}
{\small 2027 考研备考资料 · 英语二专用}\par
\vspace{12mm}
{\fontsize{31}{39}\selectfont\bfseries 英语二作文}\par
\vspace{4mm}
{\fontsize{23}{31}\selectfont\bfseries 36 类融合模板}\par
\vspace{10mm}
\rule{\linewidth}{1pt}\par
\vspace{5mm}
{\Large 26 类应用文 \quad 10 类大作文写法}\par
\vspace{6mm}
可替换模板 · 完整套用示例 · 108 条表达讲解\par
任务核对 · 数据辨析 · 逐篇文字重合检查\par
\vspace{14mm}
\begin{tabularx}{\linewidth}{@{}L{30mm}X@{}}
整理依据 & 用户提供的作文讲义、素材库、方法论笔记及真题配套资料。\\[5pt]
质量标准 & 《考研英语批改指令》第5页、第6页的英语二写作要求。\\[5pt]
词数 & 示例：应用文正文99—109词；大作文147—155词。\\[5pt]
使用原则 & 先答完整题目，再选准确表达；保留文体框架，按情境改写内容。\\
\end{tabularx}
\vfill
{\small 多来源综合新编 \quad 2026年10月25日}\par
{\small 本册的情境与数据均为练习用途；不把模板本身等同于某个确定分数。}
\clearpage
\markboth{目录}{}
\tableofcontents
\clearpage\section*{先看这两页：如何使用与如何判定质量}\addcontentsline{toc}{section}{先看这两页：如何使用与如何判定质量}
\markboth{使用与标准}{}
本册是综合改写后的英语二写作手册。A01—A26 按应用文功能与格式分类；B01—B10 按材料形式及数据关系分类。相近文体有可迁移的表达，但答复信仍须逐项回应来信，混合任务题必须完成题干的每一个动作。

\par\vspace{4pt}{\bfseries 按下面的顺序套用}\par\nopagebreak
\begin{enumerate}\item 圈出身份、对象、文体和任务动词。先在草稿纸写出任务一、任务二；若有第三项，也要落实到正文。\item 选最接近的模板，把【】中的提示全部换成英文。模板保留固定句法；示例展示同一类任务怎样写具体，不要求与模板逐字对应。\item 优先填写题目明示的信息。题目允许时可以合理补充活动细节；图表数值、比例、时间、单位则必须照材料写。\item 换题后重读整篇：主语、代词、时态、单复数、活动名称和结论要一致。删去与本题无关的句子，再统计词数。\end{enumerate}
\begin{longtable}{L{28mm}L{60mm}L{65mm}}\toprule
\textbf{占位提示} & \textbf{该填什么} & \textbf{示例}\\\midrule\endhead
动词原形 & 不带 to 的动词短语 & review the draft\\
动名词 & doing 短语 & reviewing the draft\\
完整分句 & 主语＋谓语；不自带句末标点 & the opening hours are inconvenient\\
完整句 & 包含主语、谓语及句末标点 & The proportion rose to 45\%.\\
陈述语序疑问词从句 & 嵌入问句；不再次倒装 & when the course begins\\
\bottomrule\end{longtable}
\par\vspace{4pt}{\bfseries 字数、格式和优先顺序}\par\nopagebreak
以实际题干为准。常见英语二小作文要求约100词、大作文约150词；“约”不等于固定下限。本册给的是完整示例的统计值，占位模板不标虚假的成文词数。英文缩写及连字符词通常算一词，数值算一词；具体分词方式可能产生少量差异。

正文词数不包括称呼、结束语、姓名；通知标题和落款、纪要／备忘录信息栏也另列。各示例另注明这些格式信息的词数，方便核查。考场是否需要日期、标题、邮箱或特定落款，先看题目要求。

先练书信常见任务、活动通知和静态／动态图表；然后熟悉纪要、备忘录和复合材料的结构。这是使用顺序，不是对下一次考试题型的预测。

\clearpage\par\vspace{4pt}{\bfseries 以用户给定批改指令为标准}\par\nopagebreak
以下档位来自 S09《考研英语批改指令》PDF第5、6页，是该文件整理的练习评分参考，不冒充官方评分细则。采用的是内容、语言、结构和格式的综合评价，不按复杂句数量计分，也不编造“每少10词扣1分”等规则。

\begin{longtable}{L{35mm}L{122mm}}\toprule
\textbf{小作文／10分} & \textbf{本册落实方式}\\\midrule\endhead
8—9；充分优秀可10 & 完成全部任务；语言准确自然、有变化；格式与双方关系匹配。\\
6—7 & 主要要点完整；偶有错误；组织清楚、语气适合。\\
4—5 & 多数要点覆盖；部分语言问题不妨碍整体理解。\\
2—3／1／0 & 依次对应重要遗漏或语言影响理解、基本未完成、无可评价内容。\\
\bottomrule\end{longtable}
\begin{longtable}{L{35mm}L{122mm}}\toprule
\textbf{大作文／15分} & \textbf{本册落实方式}\\\midrule\endhead
13—15 & 数据描述准确，任务完整，评论贴题且有展开；语言准确、层次清楚。\\
10—12 & 主要要点完整，语言较丰富；偶有错误，整体连贯。\\
7—9 & 多数要点覆盖；描述评论基本合理；错误不妨碍理解。\\
4—6／1—3／0 & 依次对应重要遗漏或理解偏差、严重偏题且难以表达、无可评价内容。\\
\bottomrule\end{longtable}
\par\vspace{4pt}{\bfseries 本册每篇都按这五项核对}\par\nopagebreak
\begin{enumerate}\item \textbf{任务：}请求邀请时明确邀请，请求建议时给行动和理由，要求征求意见时写出问题。\item \textbf{语言：}选择常用但准确的动词和搭配；有短句也有适量从句，不为“高级”增加错误风险。\item \textbf{组织：}每一段有任务；句子之间共享具体主题，不靠反复套用 firstly、moreover 撑起逻辑。\item \textbf{图表：}先描述可见事实，再以 may/could 等表达合理解释；静态类别、不同人群、时间趋势分别处理。\item \textbf{格式：}书信有称呼和署名；通知告示有标题和发布者；纪要记录决议与分工；备忘录突出收发信息和行动。\end{enumerate}
本册没有给模板预判“稳拿几分”。实际得分还取决于当场是否切题、信息是否完整、替换后语言是否正确及书写是否清楚。

\clearpage\section*{A01  邀请信}\addcontentsline{toc}{section}{A01  邀请信}
\markboth{A01 书信}{}
{\small 新编题：作为英语社团组织者，邀请 Smith 教授参加阅读交流会；交代活动内容、时间地点和希望对方承担的任务。}\par
{\small\textbf{行文路线：}邀请及身份 → 时间地点、流程及作用 → 确认是否出席}\par
\par\vspace{4pt}{\bfseries 可替换模板}\par\nopagebreak
\begin{writing}
\noindent Dear {\bfseries 【对方称呼】},\par
{\bfseries 【主办方】} is arranging {\bfseries 【活动名称，含冠词】}, and we would be delighted to have you as {\bfseries 【受邀身份】}.

We will meet {\bfseries 【地点＋具体时间】}. Students will first {\bfseries 【参加者要做的事，动词原形】}. We hope you can then {\bfseries 【希望对方承担的任务，动词原形】}. Your experience would help beginners {\bfseries 【对参加者的具体帮助，动词原形】}. The session should take {\bfseries 【预计时长】}, and we can adjust the speaking time to suit your schedule.

Could you let me know by {\bfseries 【回复期限】} whether you can join us?

\noindent Yours sincerely,\\Li Ming\par
\end{writing}
\par\vspace{4pt}{\bfseries 完整套用示例}\par\nopagebreak
{\small 正文 99 词；标题／称呼／落款等另 7 词}\par
\begin{example}
\noindent Dear Professor Smith,\par
Our English Society is arranging a reading workshop, and we would be delighted to have you as our guest speaker.

We will meet in Room 203 at two on Friday afternoon. Students will first discuss a short story in small groups. We hope you can then explain how to read beyond the plot and answer a few questions. Your experience would help beginners approach English fiction with greater confidence. The session should take about an hour, and we can adjust the speaking time to suit your schedule.

Could you let me know by Wednesday whether you can join us?

\noindent Yours sincerely,\\Li Ming\par
\end{example}
\par\vspace{4pt}{\bfseries 好词好句：优先记这三处}\par\nopagebreak
{\small\textbf{have you as our guest speaker} — 邀请您担任嘉宾；as 后接角色}\par
{\small\textbf{read beyond the plot} — 读懂情节背后的含义；适合阅读活动}\par
{\small\textbf{adjust ... to suit ...} — 调整……以适应……}\par
{\small\textbf{核对重点：}邀请谁做什么必须明确；to suit your schedule 只表示可协商，不要擅自修改题干已定时间。}\par
{\fontsize{8}{10}\selectfont\textbf{来源线索：}S16、S18、S19；功能与格式综合重写。}\par
\clearpage\section*{A02  推荐信：书籍／活动／地点}\addcontentsline{toc}{section}{A02  推荐信：书籍／活动／地点}
\markboth{A02 书信}{}
{\small 新编题：朋友 David 想了解本地文化，向他推荐一处周末可去的地方，说明两条理由并给出游览建议。}\par
{\small\textbf{行文路线：}推荐对象 → 两个可核实特点与需求匹配 → 实用提示}\par
\par\vspace{4pt}{\bfseries 可替换模板}\par\nopagebreak
\begin{writing}
\noindent Dear {\bfseries 【对方称呼】},\par
For {\bfseries 【对方安排】}, I would choose {\bfseries 【推荐对象】}. It seems a good match for your interest in {\bfseries 【兴趣领域】}.

One attractive feature is {\bfseries 【特点一，名词短语】}. This would allow you to {\bfseries 【具体收获，动词原形】} without {\bfseries 【可避免的困难，动名词】}. Another reason to consider it is {\bfseries 【特点二，名词短语】}, which would be especially useful when you {\bfseries 【具体场景，动词短语】}. To get the most out of the experience, I suggest {\bfseries 【准备或使用建议，动名词】}. That should leave you enough time to {\bfseries 【进一步安排，动词原形】}.

Tell me if you decide to {\bfseries 【采纳推荐的行动，动词原形】}; I would be happy to help with the arrangements.

\noindent Best wishes,\\Li Ming\par
\end{writing}
\par\vspace{4pt}{\bfseries 完整套用示例}\par\nopagebreak
{\small 正文 101 词；标题／称呼／落款等另 6 词}\par
\begin{example}
\noindent Dear David,\par
For your weekend visit, I would choose the City History Museum. It seems a good match for your interest in everyday life in China.

The displays explain how local families lived, worked and celebrated festivals. Short English descriptions make the stories easy to follow, even without much background knowledge. There is also a workshop where visitors can try traditional printing. You could therefore learn about the past while making something yourself. I suggest going early, since this would leave you enough time for both the exhibition and the workshop.

Tell me if you decide to go; I would enjoy joining you.

\noindent Best wishes,\\Li Ming\par
\end{example}
\par\vspace{4pt}{\bfseries 好词好句：优先记这三处}\par\nopagebreak
{\small\textbf{a good match for} — 与……很合适}\par
{\small\textbf{background knowledge} — 背景知识，不可数}\par
{\small\textbf{leave ... enough time for ...} — 给……留出足够时间}\par
{\small\textbf{核对重点：}推荐对象可以换成书、课程或展览，但理由也要随之改写；不要保留与新对象无关的 museum/workshop。}\par
{\fontsize{8}{10}\selectfont\textbf{来源线索：}S16、S18、S19；功能与格式综合重写。}\par
\clearpage\section*{A03  推荐信：推荐他人}\addcontentsline{toc}{section}{A03  推荐信：推荐他人}
\markboth{A03 书信}{}
{\small 新编题：向项目负责人推荐同学 Chen Yu 担任博物馆志愿讲解员，交代关系、能力及具体经历。}\par
{\small\textbf{行文路线：}推荐及关系 → 能力与可观察证据 → 明确推荐}\par
\par\vspace{4pt}{\bfseries 可替换模板}\par\nopagebreak
\begin{writing}
\noindent Dear {\bfseries 【对方称呼】},\par
May I recommend {\bfseries 【姓名】} for {\bfseries 【职位】}? We worked together on {\bfseries 【共同项目】}, so I have seen how carefully {\bfseries 【he/she】} prepares and communicates.

During that project, {\bfseries 【he/she】} took responsibility for {\bfseries 【具体任务】} and managed to {\bfseries 【可观察成果，动词原形】}. What impressed me most was {\bfseries 【表现，名词短语】}: when {\bfseries 【具体情境，完整分句】}, {\bfseries 【he/she】} {\bfseries 【应对行为，过去式动词短语】}. This experience showed both {\bfseries 【能力一】} and {\bfseries 【能力二】}, which would be useful in the role you are offering.

I believe {\bfseries 【he/she】} could make a reliable addition to your team. I can provide further details about our work if needed.

\noindent Yours sincerely,\\Li Ming\par
\end{writing}
\par\vspace{4pt}{\bfseries 完整套用示例}\par\nopagebreak
{\small 正文 103 词；标题／称呼／落款等另 7 词}\par
\begin{example}
\noindent Dear Project Coordinator,\par
May I recommend Chen Yu for the museum guide position? We worked together on a campus history project last term, so I have seen how carefully she prepares and communicates.

She can explain unfamiliar ideas in clear English and is patient with questions. During our project, she turned several difficult historical accounts into short stories that visitors could easily understand. She also checked every date against the original records. These habits would be useful when guiding groups with different levels of knowledge.

I believe she could make a reliable addition to your team. I can provide further details about our work if needed.

\noindent Yours sincerely,\\Li Ming\par
\end{example}
\par\vspace{4pt}{\bfseries 好词好句：优先记这三处}\par\nopagebreak
{\small\textbf{check ... against ...} — 对照……核查……}\par
{\small\textbf{a reliable addition to your team} — 团队中可靠的新成员}\par
{\small\textbf{if needed} — 如有需要}\par
{\small\textbf{核对重点：}she/her 应随人物一致替换；推荐人只能声称自己有依据的经历，不能只堆 excellent、outstanding。}\par
{\fontsize{8}{10}\selectfont\textbf{来源线索：}S06、S18、S29；推荐功能与证据表达综合重写。}\par
\clearpage\section*{A04  建议信}\addcontentsline{toc}{section}{A04  建议信}
\markboth{A04 书信}{}
{\small 新编题：朋友 Jack 要做英语口头报告，给出准备建议并说明这些做法有什么帮助。}\par
{\small\textbf{行文路线：}回应困难 → 可执行步骤及理由 → 后续支持}\par
\par\vspace{4pt}{\bfseries 可替换模板}\par\nopagebreak
\begin{writing}
\noindent Dear {\bfseries 【对方称呼】},\par
Your {\bfseries 【任务】} sounds interesting, and a few simple steps may make the preparation easier.

Start by {\bfseries 【第一步，动名词】}. This will help you identify {\bfseries 【需要明确的内容】} before you spend too much time on details. Next, {\bfseries 【第二步，动词原形】} so that you can {\bfseries 【具体效果，动词原形】}. You could also ask {\bfseries 【合适的人】} to {\bfseries 【反馈行为，动词原形】} and point out anything unclear. If {\bfseries 【可能困难，完整分句】}, try {\bfseries 【应对办法，动名词】}. Keep your attention on {\bfseries 【核心目标】}, since this is what the task requires.

I would be happy to {\bfseries 【愿意提供的帮助，动词原形】} when you are ready.

\noindent Best wishes,\\Li Ming\par
\end{writing}
\par\vspace{4pt}{\bfseries 完整套用示例}\par\nopagebreak
{\small 正文 101 词；标题／称呼／落款等另 6 词}\par
\begin{example}
\noindent Dear Jack,\par
Your presentation sounds interesting, and a few simple steps may make the preparation easier.

Start by choosing one question that you want your audience to understand. Collect two or three examples to answer it, instead of including every fact you find. Next, practise aloud with a timer. This will show you which parts need shortening and help you speak at a comfortable pace. You could also ask a classmate to listen and point out anything unclear. Keep the slides brief so that people can focus on what you say.

I would be happy to listen to a practice run this weekend.

\noindent Best wishes,\\Li Ming\par
\end{example}
\par\vspace{4pt}{\bfseries 好词好句：优先记这三处}\par\nopagebreak
{\small\textbf{practise aloud with a timer} — 计时大声练习}\par
{\small\textbf{at a comfortable pace} — 以舒适的速度}\par
{\small\textbf{point out anything unclear} — 指出不清楚的地方}\par
{\small\textbf{核对重点：}建议必须与来信中的困难对应；suggest 后用 doing 或 that 从句，不能写 suggest you to do。}\par
{\fontsize{8}{10}\selectfont\textbf{来源线索：}S16、S18、S19；功能与格式综合重写。}\par
\clearpage\section*{A05  计划信：提出方案并征求意见}\addcontentsline{toc}{section}{A05  计划信：提出方案并征求意见}
\markboth{A05 书信}{}
{\small 新编题：你和 Jack 准备调查校园旧物再利用情况；提出调查计划并询问他的意见。}\par
{\small\textbf{行文路线：}提出计划 → 方法、分工和时间 → 征求修改意见}\par
\par\vspace{4pt}{\bfseries 可替换模板}\par\nopagebreak
\begin{writing}
\noindent Dear {\bfseries 【对方称呼】},\par
For our survey on {\bfseries 【调查主题】}, I suggest starting with {\bfseries 【对象或地点】}. This would let us reach people with direct experience of the issue.

We could prepare {\bfseries 【数量】} short questions about {\bfseries 【调查方面一】} and {\bfseries 【调查方面二】}. I can {\bfseries 【我的任务，动词原形】}, while you {\bfseries 【对方任务，动词原形】}. After collecting the answers, we could group the main findings and discuss {\bfseries 【后续分析内容】}. {\bfseries 【拟定时间】} would work for me, but we should check when the participants are available. A brief visit beforehand could help us avoid unnecessary delays.

Does this approach sound reasonable to you? Please suggest any changes before we begin contacting people.

\noindent Best wishes,\\Li Ming\par
\end{writing}
\par\vspace{4pt}{\bfseries 完整套用示例}\par\nopagebreak
{\small 正文 102 词；标题／称呼／落款等另 6 词}\par
\begin{example}
\noindent Dear Jack,\par
For our survey on the reuse of unwanted items, I suggest starting with the student residences. This would let us reach people who regularly exchange books and household goods.

We could prepare five short questions about what students exchange and what difficulties they encounter. I can draft the questions, while you arrange interviews with a few residents. After collecting the answers, we could group the main problems and discuss possible improvements. Saturday morning would work for me, but we should check when the residents are available.

Does this approach sound reasonable to you? Please suggest any changes before we begin contacting people.

\noindent Best wishes,\\Li Ming\par
\end{example}
\par\vspace{4pt}{\bfseries 好词好句：优先记这三处}\par\nopagebreak
{\small\textbf{group the main problems} — 归类主要问题}\par
{\small\textbf{check when ... are available} — 确认……何时有空}\par
{\small\textbf{Does this approach sound reasonable to you?} — 明确征求意见}\par
{\small\textbf{核对重点：}共同计划用 could 表示提议；征求意见必须明写，不要只用期待回信收尾。}\par
{\fontsize{8}{10}\selectfont\textbf{来源线索：}S16、S18、S19；功能与格式综合重写。}\par
\clearpage\section*{A06  介绍信：介绍项目并邀请参与}\addcontentsline{toc}{section}{A06  介绍信：介绍项目并邀请参与}
\markboth{A06 书信}{}
{\small 新编题：向国际学生介绍校园故事展，并邀请他们分享家乡故事、说明参与方式。}\par
{\small\textbf{行文路线：}是什么 → 怎么参与、有什么价值 → 明确邀请与提交方式}\par
\par\vspace{4pt}{\bfseries 可替换模板}\par\nopagebreak
\begin{writing}
\noindent Dear {\bfseries 【对方称呼】},\par
Our {\bfseries 【项目名称】} will begin {\bfseries 【时间介词短语】}, and I hope you will {\bfseries 【邀请参与的动作，动词原形】}.

The project will explore {\bfseries 【主题】} through {\bfseries 【具体形式】}. Each participant will {\bfseries 【具体任务，动词原形】}, and we will provide {\bfseries 【主办方支持】} to help with the preparation. The main activity will take place {\bfseries 【地点介词短语】}, where everyone can {\bfseries 【交流或实践的动作，动词原形】}. No special experience is required; {\bfseries 【最基本条件，名词短语】} is enough. Your contribution would give others a chance to learn about {\bfseries 【相应内容】} from a different perspective.

Please send {\bfseries 【报名信息或作品】} to {\bfseries 【接收方】} by {\bfseries 【期限】}. I hope you can join us.

\noindent Best wishes,\\Li Ming\par
\end{writing}
\par\vspace{4pt}{\bfseries 完整套用示例}\par\nopagebreak
{\small 正文 103 词；标题／称呼／落款等另 7 词}\par
\begin{example}
\noindent Dear International Students,\par
Our campus storytelling exhibition will open next month, and I hope you will contribute a story from your hometown.

The exhibition will explore how ordinary places become part of people's memories. Each participant will provide one photograph and a short English description explaining its personal meaning. We will display the contributions in the student centre and hold a discussion on the opening afternoon. No special writing experience is required; a clear and honest account is enough. Your stories would give other students a chance to discover places through the people who know them.

Please send your contribution to the exhibition team by Friday.

\noindent Best wishes,\\Li Ming\par
\end{example}
\par\vspace{4pt}{\bfseries 好词好句：优先记这三处}\par\nopagebreak
{\small\textbf{contribute a story} — 提供／分享一则故事}\par
{\small\textbf{a clear and honest account} — 清楚真实的叙述}\par
{\small\textbf{No special ... is required} — 不要求专门的……}\par
{\small\textbf{核对重点：}介绍与邀请是两个任务，都要有具体内容；若题目要求邮箱或报名地点，应补全。}\par
{\fontsize{8}{10}\selectfont\textbf{来源线索：}S16、S18、S19；功能与格式综合重写。}\par
\clearpage\section*{A07  答复信：回复咨询}\addcontentsline{toc}{section}{A07  答复信：回复咨询}
\markboth{A07 书信}{}
{\small 新编题：Anna 来信询问志愿教学项目的工作、安排和准备事项，逐项回复。}\par
{\small\textbf{行文路线：}回应来信 → 按问题顺序逐项回答 → 后续联系}\par
\par\vspace{4pt}{\bfseries 可替换模板}\par\nopagebreak
\begin{writing}
\noindent Dear {\bfseries 【对方称呼】},\par
You asked about {\bfseries 【项目】}. I am glad you are considering joining us, and here are the details you need.

The main task is to {\bfseries 【工作内容，动词原形】}, with support from {\bfseries 【协助者】}. Activities take place {\bfseries 【时间地点介词短语】} and normally last {\bfseries 【时长】}. As for {\bfseries 【来信中的第二个问题】}, you will need to {\bfseries 【对应答案，动词原形】}. Before your first visit, please {\bfseries 【准备事项，动词原形】}. This will help us {\bfseries 【准备的作用，动词原形】} and make the arrangements easier for everyone. The coordinator can explain any remaining details when you arrive.

If this arrangement suits you, I can introduce you to the coordinator.

\noindent Best wishes,\\Li Ming\par
\end{writing}
\par\vspace{4pt}{\bfseries 完整套用示例}\par\nopagebreak
{\small 正文 102 词；标题／称呼／落款等另 6 词}\par
\begin{example}
\noindent Dear Anna,\par
You asked about our weekend teaching project. I am glad you are considering joining us, and here are the details you need.

Volunteers help local children practise basic English through games and short conversations. Sessions take place at the community centre on Saturday mornings, with each volunteer working alongside an experienced teacher. You do not need to prepare a full lesson alone. Before your first visit, please read the activity guide and bring one simple game idea. The coordinator will explain the children's needs and help you adapt your materials.

If this arrangement suits you, I can introduce you to the coordinator.

\noindent Best wishes,\\Li Ming\par
\end{example}
\par\vspace{4pt}{\bfseries 好词好句：优先记这三处}\par\nopagebreak
{\small\textbf{work alongside} — 与……一起工作}\par
{\small\textbf{adapt your materials} — 调整你的材料以适应需要}\par
{\small\textbf{If this arrangement suits you} — 如果这样的安排适合你}\par
{\small\textbf{核对重点：}先圈出来信里的每个问句，再检查是否一一作答；不要用笼统的活动意义代替答案。}\par
{\fontsize{8}{10}\selectfont\textbf{来源线索：}S16、S18、S19；功能与格式综合重写。}\par
\clearpage\section*{A08  交流信：表达感受并询问详情}\addcontentsline{toc}{section}{A08  交流信：表达感受并询问详情}
\markboth{A08 书信}{}
{\small 新编题：朋友 Jack 分享了家庭旅行视频；表达具体感受，并询问希望进一步了解的旅行细节。}\par
{\small\textbf{行文路线：}具体感受及细节 → 两个相关追问 → 友好收尾}\par
\par\vspace{4pt}{\bfseries 可替换模板}\par\nopagebreak
\begin{writing}
\noindent Dear {\bfseries 【对方称呼】},\par
I watched your {\bfseries 【分享内容】} and especially enjoyed the part where {\bfseries 【具体情节，完整分句】}. It made the experience feel {\bfseries 【感受形容词】}.

It was interesting to see how {\bfseries 【具体观察，完整分句】}. Your {\bfseries 【分享内容】} also made me curious about {\bfseries 【进一步想了解的方面】}. How did you {\bfseries 【具体经历，动词原形】}, and what did you enjoy most about {\bfseries 【相关活动】}? I would also love to know whether {\bfseries 【希望确认的信息，完整分句】}. These details would help me understand the experience more fully, especially since I am interested in {\bfseries 【共同兴趣】} myself.

When you have time, please share a few more stories. I would really enjoy hearing them.

\noindent Best wishes,\\Li Ming\par
\end{writing}
\par\vspace{4pt}{\bfseries 完整套用示例}\par\nopagebreak
{\small 正文 101 词；标题／称呼／落款等另 6 词}\par
\begin{example}
\noindent Dear Jack,\par
I watched your travel video yesterday and especially enjoyed the scene where your family learned to make noodles. The laughter made the experience feel warm and personal.

It was interesting to see how a simple meal brought you together with local people. Your video also made me curious about the rest of the journey. How did you find the cooking workshop, and which place surprised you most? I would also love to know whether travelling with your parents changed the way you planned each day.

When you have time, please share a few more stories. I would really enjoy hearing them.

\noindent Best wishes,\\Li Ming\par
\end{example}
\par\vspace{4pt}{\bfseries 好词好句：优先记这三处}\par\nopagebreak
{\small\textbf{made me curious about} — 使我想进一步了解……}\par
{\small\textbf{which place surprised you most} — 哪个地方最让你意外}\par
{\small\textbf{warm and personal} — 温暖且有亲近感}\par
{\small\textbf{核对重点：}不是写一封正式咨询公函；询问必须落到旅行或视频细节，不能只说 tell me more。}\par
{\fontsize{8}{10}\selectfont\textbf{来源线索：}S19、S36、S40；交流任务据所给题干改编，措辞与情节新编。}\par
\clearpage\section*{A09  感谢信}\addcontentsline{toc}{section}{A09  感谢信}
\markboth{A09 书信}{}
{\small 新编题：感谢 Smith 教授在你准备学术报告时给予的具体帮助，说明帮助带来的变化。}\par
{\small\textbf{行文路线：}感谢具体事项 → 帮助及其效果 → 延伸收获、再次致谢}\par
\par\vspace{4pt}{\bfseries 可替换模板}\par\nopagebreak
\begin{writing}
\noindent Dear {\bfseries 【对方称呼】},\par
Now that {\bfseries 【受助事项完成，完整分句】}, I want to thank you for {\bfseries 【具体帮助，动名词】}. Your advice made a real difference to {\bfseries 【改进方面，名词或动名词】}.

When I struggled with {\bfseries 【遇到的困难】}, you helped me {\bfseries 【帮助一，动词原形】}. You also showed me how to {\bfseries 【帮助二，动词原形】}, which made the next stage much easier. As a result, I was able to {\bfseries 【具体成果，动词原形】} with greater confidence. More than the result itself, I value the {\bfseries 【长期收获】} that your support helped me develop.

I will use what I learned in future work. Thank you again for your time and patience.

\noindent Yours sincerely,\\Li Ming\par
\end{writing}
\par\vspace{4pt}{\bfseries 完整套用示例}\par\nopagebreak
{\small 正文 108 词；标题／称呼／落款等另 7 词}\par
\begin{example}
\noindent Dear Professor Smith,\par
Now that my presentation is over, I want to thank you for helping me prepare it. Your advice made a real difference to how I explained my ideas.

When my first draft was too broad, you helped me choose a clear question and remove unrelated material. You also showed me how to use examples instead of reading directly from my notes. As a result, I felt more confident and could respond calmly to the audience's questions. More than the presentation itself, I value the careful way of thinking that you encouraged.

I will use what I learned in future work. Thank you again for your time and patience.

\noindent Yours sincerely,\\Li Ming\par
\end{example}
\par\vspace{4pt}{\bfseries 好词好句：优先记这三处}\par\nopagebreak
{\small\textbf{make a real difference to} — 对……有切实影响}\par
{\small\textbf{respond calmly to} — 冷静回应……}\par
{\small\textbf{your time and patience} — 您付出的时间和耐心}\par
{\small\textbf{核对重点：}感谢信需要“你做了什么＋这对我有什么用”；不用频繁重复 heartfelt/sincere 来凑词。}\par
{\fontsize{8}{10}\selectfont\textbf{来源线索：}S16、S18、S19；功能与格式综合重写。}\par
\clearpage\section*{A10  道歉信}\addcontentsline{toc}{section}{A10  道歉信}
\markboth{A10 书信}{}
{\small 新编题：借用朋友 Alex 的书后不慎损坏，致歉、说明经过并提出具体补救方案。}\par
{\small\textbf{行文路线：}具体道歉 → 经过、承担责任、可行补救 → 再次致歉}\par
\par\vspace{4pt}{\bfseries 可替换模板}\par\nopagebreak
\begin{writing}
\noindent Dear {\bfseries 【对方称呼】},\par
I owe you an apology for {\bfseries 【过失，动名词】}. Unfortunately, {\bfseries 【事件经过，完整分句】}, which caused {\bfseries 【对方受到的影响】}.

I should have {\bfseries 【本应采取的行动，过去分词短语】}, and I take responsibility for the inconvenience. Since you need {\bfseries 【对方实际需要】}, I have arranged to {\bfseries 【补救动作，动词原形】} by {\bfseries 【完成期限】}. Until then, I can {\bfseries 【临时解决办法，动词原形】} so that you can continue with your plans. If this arrangement causes any difficulty, please tell me and I will find another solution.

I am sorry for the trouble this has caused. I will be more careful about {\bfseries 【需注意的事项】} in future.

\noindent Best wishes,\\Li Ming\par
\end{writing}
\par\vspace{4pt}{\bfseries 完整套用示例}\par\nopagebreak
{\small 正文 104 词；标题／称呼／落款等另 6 词}\par
\begin{example}
\noindent Dear Alex,\par
I owe you an apology for damaging the book you lent me. I accidentally spilled water on it yesterday, and several pages are now stained.

I should have kept it away from my desk while I was drinking. I understand that you need the book for your course, so I have ordered a replacement of the same edition. It should arrive on Thursday. Until then, you are welcome to use my copy, which covers the same material. If this arrangement causes any difficulty, please tell me and I will find another solution.

I am sorry for being careless and for disrupting your study plans.

\noindent Best wishes,\\Li Ming\par
\end{example}
\par\vspace{4pt}{\bfseries 好词好句：优先记这三处}\par\nopagebreak
{\small\textbf{owe you an apology for} — 应为……向你道歉}\par
{\small\textbf{a replacement of the same edition} — 同一版本的替换本}\par
{\small\textbf{disrupt your study plans} — 打乱你的学习计划}\par
{\small\textbf{核对重点：}解释原因不能替代承担责任；不要只说以后注意而不给对方可执行的解决办法。}\par
{\fontsize{8}{10}\selectfont\textbf{来源线索：}S16、S18、S19；功能与格式综合重写。}\par
\clearpage\section*{A11  投诉信}\addcontentsline{toc}{section}{A11  投诉信}
\markboth{A11 书信}{}
{\small 新编题：从商店购买的台灯无法正常使用；说明问题、影响和要求，语气坚定而礼貌。}\par
{\small\textbf{行文路线：}交易及诉求 → 问题、影响与证据 → 处理方式和期限}\par
\par\vspace{4pt}{\bfseries 可替换模板}\par\nopagebreak
\begin{writing}
\noindent Dear {\bfseries 【对方称呼】},\par
The {\bfseries 【物品】} I bought from your store {\bfseries 【购买时间】} developed a problem {\bfseries 【问题出现时间】}. I am contacting you to request {\bfseries 【合理诉求】}.

Although I followed the instructions, {\bfseries 【具体故障，完整分句】}. This has made it difficult for me to {\bfseries 【正常用途，动词原形】}. I have already tried {\bfseries 【排查措施，动名词】}, but the problem remains. I can provide the receipt and {\bfseries 【其他证据】} if they would help you assess the situation. I would prefer {\bfseries 【具体处理方案】}, provided it can be arranged promptly.

Please explain the return procedure and tell me when I can expect a solution. A reply by {\bfseries 【期限】} would be appreciated.

\noindent Yours sincerely,\\Li Ming\par
\end{writing}
\par\vspace{4pt}{\bfseries 完整套用示例}\par\nopagebreak
{\small 正文 104 词；标题／称呼／落款等另 8 词}\par
\begin{example}
\noindent Dear Customer Service Team,\par
The desk lamp I bought from your store on Monday stopped working after two evenings. I am contacting you to request a replacement.

Although I followed the instructions, the light now switches off whenever I adjust its angle. This makes it unsuitable for reading and has interrupted my evening study. I have checked the power supply, so the problem appears to be with the lamp itself. The receipt and photographs of the item are available if you need them.

Please let me know how to return it and when a replacement can be provided. I would appreciate a reply within the next few days.

\noindent Yours sincerely,\\Li Ming\par
\end{example}
\par\vspace{4pt}{\bfseries 好词好句：优先记这三处}\par\nopagebreak
{\small\textbf{request a replacement} — 请求换货}\par
{\small\textbf{appears to be with ...} — 看来出在……上}\par
{\small\textbf{The receipt ... is available} — 可提供收据；注意主谓一致}\par
{\small\textbf{核对重点：}事实与要求应清楚；没有题干依据，不写人身攻击、威胁或夸大的损失。}\par
{\fontsize{8}{10}\selectfont\textbf{来源线索：}S16、S18、S19；功能与格式综合重写。}\par
\clearpage\section*{A12  咨询信}\addcontentsline{toc}{section}{A12  咨询信}
\markboth{A12 书信}{}
{\small 新编题：准备参加暑期语言课程，向负责人询问课程安排、费用和报名条件。}\par
{\small\textbf{行文路线：}背景与目的 → 三个相关问题 → 说明信息用途并致谢}\par
\par\vspace{4pt}{\bfseries 可替换模板}\par\nopagebreak
\begin{writing}
\noindent Dear {\bfseries 【对方称呼】},\par
I am considering {\bfseries 【项目】} and would like to clarify a few points before applying. My main aim is to {\bfseries 【个人目标，动词原形】}.

Could you tell me {\bfseries 【问题一，陈述语序的疑问词从句】}? I would also like to know whether {\bfseries 【问题二，完整分句】}. As I will be {\bfseries 【相关背景，动名词或状态】}, information about {\bfseries 【额外信息】} would be useful. Finally, please explain {\bfseries 【问题三，名词或疑问词从句】} and let me know if any supporting documents are needed. I hope to make the necessary arrangements {\bfseries 【时间介词短语】}.

These details would help me plan ahead. Thank you for taking the time to answer my questions.

\noindent Yours sincerely,\\Li Ming\par
\end{writing}
\par\vspace{4pt}{\bfseries 完整套用示例}\par\nopagebreak
{\small 正文 107 词；标题／称呼／落款等另 7 词}\par
\begin{example}
\noindent Dear Course Coordinator,\par
I am considering your summer language course and would like to clarify a few points before applying. My main aim is to become more confident in academic discussions.

Could you tell me how many hours of speaking practice are included each week? I would also like to know whether the fee covers learning materials or whether these must be purchased separately. As I will be travelling from another city, information about nearby accommodation would be useful. Finally, is an entrance test required, and when should applicants submit their results?

These details would help me plan my stay. Thank you for taking the time to answer my questions.

\noindent Yours sincerely,\\Li Ming\par
\end{example}
\par\vspace{4pt}{\bfseries 好词好句：优先记这三处}\par\nopagebreak
{\small\textbf{clarify a few points} — 澄清几个问题}\par
{\small\textbf{be purchased separately} — 需要另行购买}\par
{\small\textbf{submit their results} — 提交成绩／结果}\par
{\small\textbf{核对重点：}Could you tell me 后用陈述语序：how many hours are included；不是 are how many hours。题目已有信息不重复问。}\par
{\fontsize{8}{10}\selectfont\textbf{来源线索：}S16、S18、S19；功能与格式综合重写。}\par
\clearpage\section*{A13  请求信：请求帮助}\addcontentsline{toc}{section}{A13  请求信：请求帮助}
\markboth{A13 书信}{}
{\small 新编题：希望 Smith 教授审阅调查问卷，说明请求内容、困难和时间安排。}\par
{\small\textbf{行文路线：}明确请求 → 为什么需要帮助、准备了什么 → 尊重对方时间}\par
\par\vspace{4pt}{\bfseries 可替换模板}\par\nopagebreak
\begin{writing}
\noindent Dear {\bfseries 【对方称呼】},\par
Could you spare some time to {\bfseries 【请求事项，动词原形】}? I would value your advice before {\bfseries 【相关期限或事件】}.

I have already {\bfseries 【已完成工作，过去分词短语】}, but I am unsure about {\bfseries 【具体困难】}. Your experience in {\bfseries 【专长领域】} would help me {\bfseries 【希望达成的效果，动词原形】}. I can send you {\bfseries 【材料】} today, together with a short explanation of {\bfseries 【背景】}. Comments on {\bfseries 【请求范围】} would be enough, so there is no need to review the whole project. I would be happy to adjust the timing if necessary.

Please let me know whether this would fit your schedule. Thank you for considering the request.

\noindent Yours sincerely,\\Li Ming\par
\end{writing}
\par\vspace{4pt}{\bfseries 完整套用示例}\par\nopagebreak
{\small 正文 109 词；标题／称呼／落款等另 7 词}\par
\begin{example}
\noindent Dear Professor Smith,\par
Could you spare some time to review the questionnaire for our student reading survey? I would value your advice before we distribute it next week.

Our team has drafted eight questions, but we are unsure whether two of them might lead respondents towards a particular answer. Your experience in survey design would help us make the wording more balanced. I can send you the draft today, together with a short explanation of our research aim. Comments on those two questions would be enough, so a detailed review of the whole project is unnecessary.

Please let me know whether this would fit your schedule. Thank you for considering the request.

\noindent Yours sincerely,\\Li Ming\par
\end{example}
\par\vspace{4pt}{\bfseries 好词好句：优先记这三处}\par\nopagebreak
{\small\textbf{spare some time to do} — 抽时间做……}\par
{\small\textbf{lead respondents towards a particular answer} — 诱导受访者选择特定答案}\par
{\small\textbf{fit your schedule} — 符合您的时间安排}\par
{\small\textbf{核对重点：}求助要说明具体任务和截止时间；不能把对方同意当作既成事实。}\par
{\fontsize{8}{10}\selectfont\textbf{来源线索：}S16、S18、S19；功能与格式综合重写。}\par
\clearpage\section*{A14  申请信：项目／志愿者}\addcontentsline{toc}{section}{A14  申请信：项目／志愿者}
\markboth{A14 书信}{}
{\small 新编题：申请担任校园文化节志愿者；说明动机、相关经历和可服务时间。}\par
{\small\textbf{行文路线：}申请及动机 → 经历证据、能力和时间 → 面谈意愿}\par
\par\vspace{4pt}{\bfseries 可替换模板}\par\nopagebreak
\begin{writing}
\noindent Dear {\bfseries 【对方称呼】},\par
Please consider my application for {\bfseries 【项目或志愿岗位】}. I would enjoy {\bfseries 【具体服务，动名词】} and contributing to {\bfseries 【服务目标】}.

During {\bfseries 【相关经历】}, I helped {\bfseries 【具体工作，动词原形】} and learned to {\bfseries 【相关能力，动词原形】}. When {\bfseries 【实际困难，完整分句】}, I {\bfseries 【应对行为，过去式动词短语】}. This experience would be useful when carrying out the duties you described. I can also {\bfseries 【其他能力，动词原形】} and am willing to complete any training required. I will be available {\bfseries 【服务时间介词短语】}.

I would welcome a brief interview to discuss how I could contribute. Thank you for reviewing my application.

\noindent Yours sincerely,\\Li Ming\par
\end{writing}
\par\vspace{4pt}{\bfseries 完整套用示例}\par\nopagebreak
{\small 正文 104 词；标题／称呼／落款等另 7 词}\par
\begin{example}
\noindent Dear Festival Coordinator,\par
Please consider my application to volunteer at the campus cultural festival. I would enjoy helping visitors explore the exhibition and feel welcome on campus.

Last semester, I helped organise a student art show, where I answered questions and guided visitors between rooms. That experience taught me to give clear directions and remain patient when plans changed. I can communicate comfortably in English and am willing to learn the background of each display before the festival begins. I will be available throughout both days of the event.

I would welcome a brief interview to discuss how I could contribute. Thank you for reviewing my application.

\noindent Yours sincerely,\\Li Ming\par
\end{example}
\par\vspace{4pt}{\bfseries 好词好句：优先记这三处}\par\nopagebreak
{\small\textbf{Please consider my application to ...} — 请考虑我申请……}\par
{\small\textbf{remain patient when plans changed} — 安排变化时仍保持耐心}\par
{\small\textbf{be available throughout ...} — 在……全程有空}\par
{\small\textbf{核对重点：}按申请要求选经历，不编造题目相冲突的身份；如申请奖学金或交换项目，要改成对应资格。}\par
{\fontsize{8}{10}\selectfont\textbf{来源线索：}S16、S18、S19；功能与格式综合重写。}\par
\clearpage\section*{A15  求职信}\addcontentsline{toc}{section}{A15  求职信}
\markboth{A15 书信}{}
{\small 新编题：申请校报编辑助理；写明申请岗位、与工作匹配的技能和可到岗时间。}\par
{\small\textbf{行文路线：}申请岗位 → 经历与职责匹配 → 到岗时间和面试}\par
\par\vspace{4pt}{\bfseries 可替换模板}\par\nopagebreak
\begin{writing}
\noindent Dear {\bfseries 【对方称呼】},\par
Your advertisement for {\bfseries 【岗位】} caught my attention because the role involves {\bfseries 【相关工作，动名词】}. I would like to apply for the position.

As {\bfseries 【身份，含冠词】}, I have gained experience in {\bfseries 【相关领域】} through {\bfseries 【实际经历】}. My responsibilities included {\bfseries 【具体职责，动名词】}, which taught me to {\bfseries 【匹配能力，动词原形】}. I also {\bfseries 【另一项技能，动词短语】}, so I could support the team with {\bfseries 【岗位需要】}. I can start {\bfseries 【到岗时间】} and would be willing to discuss a suitable schedule.

My resume provides further details of my experience. I would appreciate an opportunity to discuss the role in an interview.

\noindent Yours sincerely,\\Li Ming\par
\end{writing}
\par\vspace{4pt}{\bfseries 完整套用示例}\par\nopagebreak
{\small 正文 105 词；标题／称呼／落款等另 7 词}\par
\begin{example}
\noindent Dear Hiring Manager,\par
Your advertisement for an editorial assistant caught my attention because the role combines writing with careful fact checking. I would like to apply for the position.

As an English major, I have edited articles for our student newsletter and worked with writers to improve clarity. I also maintain a simple checklist for names, dates and sources, which helps me catch mistakes before publication. These experiences would allow me to support your editors with both language and accuracy. I can start after the current term ends.

My resume provides further details of my experience. I would appreciate an opportunity to discuss the role in an interview.

\noindent Yours sincerely,\\Li Ming\par
\end{example}
\par\vspace{4pt}{\bfseries 好词好句：优先记这三处}\par\nopagebreak
{\small\textbf{careful fact checking} — 仔细核查事实}\par
{\small\textbf{improve clarity} — 提升表达清晰度}\par
{\small\textbf{before publication} — 发表之前}\par
{\small\textbf{核对重点：}简历若未要求随附，可写 provides details；考试允许合理补细节，但不能用真实姓名或违反题干身份。}\par
{\fontsize{8}{10}\selectfont\textbf{来源线索：}S05、S18、S29；岗位申请功能综合重写。}\par
\clearpage\section*{A16  祝贺信}\addcontentsline{toc}{section}{A16  祝贺信}
\markboth{A16 书信}{}
{\small 新编题：朋友 Anna 获得演讲比赛奖项，向她祝贺并谈及你观察到的努力。}\par
{\small\textbf{行文路线：}直接祝贺 → 努力细节与成绩相连 → 祝愿或庆祝}\par
\par\vspace{4pt}{\bfseries 可替换模板}\par\nopagebreak
\begin{writing}
\noindent Dear {\bfseries 【对方称呼】},\par
What wonderful news that {\bfseries 【获奖或喜讯，完整分句】}! You have every reason to feel proud, and I am very happy for you.

I still remember {\bfseries 【见证过的努力，名词性从句】}. {\bfseries 【另一项具体努力，完整句】} The award recognises that effort as well as your ability to {\bfseries 【相关能力，动词原形】}. I hope the experience encourages you to {\bfseries 【未来祝愿，动词原形】}.

When you are free, let us celebrate over lunch. I would love to hear which part of {\bfseries 【对应活动阶段】} you enjoyed most.

\noindent Best wishes,\\Li Ming\par
\end{writing}
\par\vspace{4pt}{\bfseries 完整套用示例}\par\nopagebreak
{\small 正文 106 词；标题／称呼／落款等另 6 词}\par
\begin{example}
\noindent Dear Anna,\par
What wonderful news that you won the university speaking competition! You have every reason to feel proud, and I am very happy for you.

I still remember how carefully you revised your speech and practised the opening until it sounded natural. You also listened seriously to other people's suggestions, even when making changes took extra time. The award recognises that effort as well as your ability to connect with an audience. I hope the experience encourages you to take on further speaking opportunities.

When you are free, let us celebrate over lunch. I would love to hear which part of the final round you enjoyed most.

\noindent Best wishes,\\Li Ming\par
\end{example}
\par\vspace{4pt}{\bfseries 好词好句：优先记这三处}\par\nopagebreak
{\small\textbf{have every reason to feel proud} — 完全有理由为自己自豪}\par
{\small\textbf{connect with an audience} — 与听众建立交流}\par
{\small\textbf{recognise that effort} — 肯定这份努力}\par
{\small\textbf{核对重点：}祝贺 on an achievement；不要写 congratulate you for。对方的努力应与所获成绩相关。}\par
{\fontsize{8}{10}\selectfont\textbf{来源线索：}S16、S18、S19；功能与格式综合重写。}\par
\clearpage\section*{A17  欢迎信}\addcontentsline{toc}{section}{A17  欢迎信}
\markboth{A17 书信}{}
{\small 新编题：以校园伙伴身份欢迎新来的国际学生，介绍报到安排并提供帮助。}\par
{\small\textbf{行文路线：}欢迎及身份 → 报到、材料和适应安排 → 联系帮助}\par
\par\vspace{4pt}{\bfseries 可替换模板}\par\nopagebreak
\begin{writing}
\noindent Dear {\bfseries 【对方称呼】},\par
Welcome to {\bfseries 【欢迎加入的单位】}! As {\bfseries 【你的身份】}, I look forward to helping you settle into campus life {\bfseries 【协助时段】}.

We will meet {\bfseries 【见面地点和时间】}. {\bfseries 【具体安排，完整句】} Please bring {\bfseries 【要准备的材料】}, as the staff will need them to {\bfseries 【材料用途，动词原形】}. There will also be {\bfseries 【后续活动，含冠词】} where you can meet other new students and ask questions about daily life.

If {\bfseries 【可能需要协调的事项】} changes, send me a message. I hope your first days here will be comfortable and enjoyable.

\noindent Best wishes,\\Li Ming\par
\end{writing}
\par\vspace{4pt}{\bfseries 完整套用示例}\par\nopagebreak
{\small 正文 105 词；标题／称呼／落款等另 7 词}\par
\begin{example}
\noindent Dear New Students,\par
Welcome to our university! As your student buddy, I look forward to helping you settle into campus life during your first week.

We will meet at the international office at nine on Monday morning. After registration, I can show you the library, dining hall and student service centre. Please bring your passport and admission letter, as the staff will need them to complete the paperwork. There will also be an informal gathering where you can meet other new students and ask questions about daily life.

If your arrival time changes, send me a message. I hope your first days here will be comfortable and enjoyable.

\noindent Best wishes,\\Li Ming\par
\end{example}
\par\vspace{4pt}{\bfseries 好词好句：优先记这三处}\par\nopagebreak
{\small\textbf{settle into campus life} — 适应校园生活}\par
{\small\textbf{complete the paperwork} — 完成手续}\par
{\small\textbf{an informal gathering} — 轻松的聚会}\par
{\small\textbf{核对重点：}welcome to 后接地点或机构；welcome you to 后才可接动词。欢迎信要有帮助新人落地的信息。}\par
{\fontsize{8}{10}\selectfont\textbf{来源线索：}S16、S18、S19；功能与格式综合重写。}\par
\clearpage\section*{A18  辞职信}\addcontentsline{toc}{section}{A18  辞职信}
\markboth{A18 书信}{}
{\small 新编题：兼职编辑工作与毕业研究冲突，向 Wang 先生说明辞职决定、原因和交接办法。}\par
{\small\textbf{行文路线：}明确决定与离职日期 → 合理原因、交接安排 → 感谢}\par
\par\vspace{4pt}{\bfseries 可替换模板}\par\nopagebreak
\begin{writing}
\noindent Dear {\bfseries 【对方称呼】},\par
After careful consideration, I have decided to leave my position as {\bfseries 【当前岗位，含冠词】} {\bfseries 【离职时间】}. I am grateful for the guidance you have given me.

{\bfseries 【辞职原因，完整句】} Rather than let unfinished work affect the team, I would like to arrange a smooth handover. I will {\bfseries 【交接措施，动词原形】}. I can also {\bfseries 【补充协助，动词原形】} to whoever takes over.

Thank you for your understanding. Working with the team has been a valuable learning experience.

\noindent Yours sincerely,\\Li Ming\par
\end{writing}
\par\vspace{4pt}{\bfseries 完整套用示例}\par\nopagebreak
{\small 正文 106 词；标题／称呼／落款等另 7 词}\par
\begin{example}
\noindent Dear Mr. Wang,\par
After careful consideration, I have decided to leave my position as a part-time editor at the end of this month. I am grateful for the guidance you have given me.

My final research project now requires more time than I expected, and I cannot give both responsibilities the attention they deserve. Rather than let unfinished work affect the team, I would like to arrange a smooth handover. I will complete my current articles and prepare notes on all outstanding tasks. I can also explain our editing process to whoever takes over.

Thank you for your understanding. Working with the team has been a valuable learning experience.

\noindent Yours sincerely,\\Li Ming\par
\end{example}
\par\vspace{4pt}{\bfseries 好词好句：优先记这三处}\par\nopagebreak
{\small\textbf{arrange a smooth handover} — 安排顺利交接}\par
{\small\textbf{outstanding tasks} — 尚未完成的工作}\par
{\small\textbf{whoever takes over} — 接手的人}\par
{\small\textbf{核对重点：}不要只说抱歉而不说辞职；outstanding tasks 在此指未完成事项，不是优秀任务。}\par
{\fontsize{8}{10}\selectfont\textbf{来源线索：}S05、S18、S31；辞职功能与交接思路综合重写。}\par
\clearpage\section*{A19  婉拒信}\addcontentsline{toc}{section}{A19  婉拒信}
\markboth{A19 书信}{}
{\small 新编题：感谢朋友邀请参加周六活动，解释不能出席的原因并提出其他参与方式。}\par
{\small\textbf{行文路线：}感谢并婉拒 → 合理原因与替代办法 → 祝愿及确认}\par
\par\vspace{4pt}{\bfseries 可替换模板}\par\nopagebreak
\begin{writing}
\noindent Dear {\bfseries 【对方称呼】},\par
Thank you for inviting me to {\bfseries 【对方邀请的活动，动词原形】}. I would have enjoyed taking part, but I cannot attend {\bfseries 【无法出席的时间】}.

{\bfseries 【无法出席的原因，完整句】} I am sorry to miss {\bfseries 【对活动价值的认可】}. If it would help, I could {\bfseries 【两个替代协助方案，动词原形】}. That way, I could still support the team without disrupting the arrangements.

I hope the day goes well. Please tell me whether either of these alternatives would be useful.

\noindent Best wishes,\\Li Ming\par
\end{writing}
\par\vspace{4pt}{\bfseries 完整套用示例}\par\nopagebreak
{\small 正文 105 词；标题／称呼／落款等另 6 词}\par
\begin{example}
\noindent Dear Alex,\par
Thank you for inviting me to help at the community book fair. I would have enjoyed taking part, but I cannot attend this Saturday.

I have already agreed to accompany my grandfather to a medical appointment that morning, and it cannot easily be rearranged. I am sorry to miss an event that brings books to more families. If it would help, I could sort the donated books on Friday evening or prepare labels before the fair. That way, I could still support the team without disrupting the arrangements.

I hope the day goes well. Please tell me whether either of these alternatives would be useful.

\noindent Best wishes,\\Li Ming\par
\end{example}
\par\vspace{4pt}{\bfseries 好词好句：优先记这三处}\par\nopagebreak
{\small\textbf{cannot easily be rearranged} — 不容易改期}\par
{\small\textbf{without disrupting the arrangements} — 不打乱已有安排}\par
{\small\textbf{either of these alternatives} — 这两个替代方案中的任何一个}\par
{\small\textbf{核对重点：}未承诺参加时主要是婉拒；已答应却失约时改用 A10 道歉逻辑，不虚构很严重的事故作借口。}\par
{\fontsize{8}{10}\selectfont\textbf{来源线索：}S16、S18、S19；功能与格式综合重写。}\par
\clearpage\section*{A20  活动通知}\addcontentsline{toc}{section}{A20  活动通知}
\markboth{A20 通知}{}
{\small 新编题：学生会发布旧物修理工作坊通知，介绍活动、参加要求及报名办法。}\par
{\small\textbf{行文路线：}活动及对象 → 内容、要求和保障 → 报名行动}\par
\par\vspace{4pt}{\bfseries 可替换模板}\par\nopagebreak
\begin{writing}
{\centering NOTICE\\{\bfseries 【具体活动或招募标题】}\par}
{\bfseries 【活动名称，含冠词】} will take place {\bfseries 【时间地点介词短语】}. All {\bfseries 【目标人群】} interested in {\bfseries 【相关兴趣，动名词】} are welcome.

The session will begin with {\bfseries 【环节一】}, followed by {\bfseries 【环节二】}. Participants will be able to {\bfseries 【实践机会，动词原形】} with help from {\bfseries 【指导者】}. Please bring {\bfseries 【所需用品】} and arrive {\bfseries 【提前量】} before the starting time. {\bfseries 【主办方提供的用品，复数】} will be available at the venue. To keep the activity running smoothly, everyone should {\bfseries 【具体要求，动词原形】}.

To reserve a place, register {\bfseries 【报名方式及期限】}. Places are limited, so early registration is recommended. Come and try something useful.

\noindent {\bfseries 【题目要求的署名或发布单位】}\par
\end{writing}
\par\vspace{4pt}{\bfseries 完整套用示例}\par\nopagebreak
{\small 正文 106 词；标题／称呼／落款等另 5 词}\par
\begin{example}
{\centering NOTICE\\Repair Workshop\par}
A repair workshop will take place in the student centre from two to four on Friday afternoon. All students interested in giving old items a second life are welcome.

Experienced volunteers will demonstrate basic repairs and guide participants through simple tasks. Please bring one small item, such as a bag or a piece of clothing, that needs attention. Electrical appliances cannot be accepted because specialist equipment is unavailable. Tools and basic materials will be provided at the venue.

To reserve a place, register at the student service desk by Wednesday. Places are limited, so early registration is recommended. Come and discover what you can repair yourself.

\noindent Students' Union\par
\end{example}
\par\vspace{4pt}{\bfseries 好词好句：优先记这三处}\par\nopagebreak
{\small\textbf{give old items a second life} — 让旧物重新发挥用途}\par
{\small\textbf{at the venue} — 在活动场地}\par
{\small\textbf{reserve a place} — 预约名额}\par
{\small\textbf{核对重点：}通知有标题和发布者，不写 Dear 或 Yours。落款与日期服从题目，不必擅加个人身份。}\par
{\fontsize{8}{10}\selectfont\textbf{来源线索：}S16、S18、S19、S31；通知结构与活动信息综合重写。}\par
\clearpage\section*{A21  招募通知}\addcontentsline{toc}{section}{A21  招募通知}
\markboth{A21 通知}{}
{\small 新编题：为口述历史项目招募四名学生志愿者，写清职责、要求、培训与申请方式。}\par
{\small\textbf{行文路线：}招募岗位 → 职责、条件和培训 → 如何申请}\par
\par\vspace{4pt}{\bfseries 可替换模板}\par\nopagebreak
\begin{writing}
{\centering NOTICE\\{\bfseries 【具体活动或招募标题】}\par}
{\bfseries 【人数】} student volunteers are needed for {\bfseries 【项目】} {\bfseries 【时间介词短语】}. The team will work with {\bfseries 【服务对象】} to {\bfseries 【目标，动词原形】}.

Volunteers will help {\bfseries 【职责一，动词原形】} and {\bfseries 【职责二，动词原形】}. Applicants should be able to {\bfseries 【能力要求，动词原形】} and be available {\bfseries 【时间要求】}. Experience in {\bfseries 【相关领域】} would be helpful, but training will be provided. Everyone must {\bfseries 【特别要求，动词原形】} when working with the team. These arrangements will help ensure that {\bfseries 【工作标准，完整分句】}.

Please submit {\bfseries 【申请材料】} to {\bfseries 【接收方】} by {\bfseries 【期限】}. Include your available times and any relevant experience. Selected applicants will be contacted for a brief meeting.

\noindent {\bfseries 【题目要求的署名或发布单位】}\par
\end{writing}
\par\vspace{4pt}{\bfseries 完整套用示例}\par\nopagebreak
{\small 正文 107 词；标题／称呼／落款等另 5 词}\par
\begin{example}
{\centering NOTICE\\Volunteers Wanted\par}
Four student volunteers are needed for a community oral history project next month. The team will record residents' memories and prepare short written summaries.

Volunteers will help arrange interviews, check recordings and organise photographs. Applicants should be patient listeners, write clearly in English and be available on two Saturday mornings. Previous interviewing experience would be helpful, but training will be provided. All participants must respect residents' privacy and ask permission before recording or sharing their stories.

Please submit a short statement of interest to the student service desk by Friday. Include your available times and any relevant experience. Selected applicants will be contacted for a brief meeting.

\noindent Project Team\par
\end{example}
\par\vspace{4pt}{\bfseries 好词好句：优先记这三处}\par\nopagebreak
{\small\textbf{a statement of interest} — 简短的申请意向说明}\par
{\small\textbf{relevant experience} — 相关经验}\par
{\small\textbf{Selected applicants will be contacted} — 将联系入选的申请者}\par
{\small\textbf{核对重点：}duties 和 requirements 必须分别写到；不能只介绍活动意义或罗列性格词。}\par
{\fontsize{8}{10}\selectfont\textbf{来源线索：}S16、S18、S19、S29；招募任务综合重写。}\par
\clearpage\section*{A22  寻物启事}\addcontentsline{toc}{section}{A22  寻物启事}
\markboth{A22 告示}{}
{\small 新编题：你在图书馆丢失蓝色笔记本，发布启事，给出特征、可能遗失地点和联系方式。}\par
{\small\textbf{行文路线：}丢失物品、时间地点 → 特征及用途 → 联系与致谢}\par
\par\vspace{4pt}{\bfseries 可替换模板}\par\nopagebreak
\begin{writing}
{\centering LOST\\{\bfseries 【遗失物品】}\par}
I am looking for {\bfseries 【丢失物品，含冠词】} that I lost {\bfseries 【丢失地点及时间】}. I last used it {\bfseries 【最后使用地点】}.

{\bfseries 【外观特征，完整句】} It contains {\bfseries 【物品内容或用途】}, so it is particularly important to me. {\bfseries 【辨认线索，完整句】} If you have found it or handed it to a member of staff, please let me know where it can be collected.

You can leave a message for Li Ming {\bfseries 【联系方式或联系地点】}. Any help would be greatly appreciated.

\noindent {\bfseries 【题目要求的署名或发布单位】}\par
\end{writing}
\par\vspace{4pt}{\bfseries 完整套用示例}\par\nopagebreak
{\small 正文 107 词；标题／称呼／落款等另 6 词}\par
\begin{example}
{\centering LOST\\A Blue Notebook\par}
I am looking for a blue notebook that I lost in the university library on Tuesday afternoon. I last used it at a desk near the second-floor windows.

The notebook has a plain cover and a small white label on the back. It contains handwritten notes for my final research project, so it is particularly important to me. My name is written inside the front cover. If you have found it or handed it to a member of staff, please let me know where it can be collected.

You can leave a message for Li Ming at the library service desk. Any help would be greatly appreciated.

\noindent Li Ming\par
\end{example}
\par\vspace{4pt}{\bfseries 好词好句：优先记这三处}\par\nopagebreak
{\small\textbf{a plain cover} — 素色封面}\par
{\small\textbf{hand it to a member of staff} — 把它交给工作人员}\par
{\small\textbf{where it can be collected} — 到哪里领取}\par
{\small\textbf{核对重点：}不必虚构高额酬金；考试中的姓名与联系方式须服从题目。此处联系地点为新编情境。}\par
{\fontsize{8}{10}\selectfont\textbf{来源线索：}S05、S18、S19；告示信息功能综合重写。}\par
\clearpage\section*{A23  招领启事}\addcontentsline{toc}{section}{A23  招领启事}
\markboth{A23 告示}{}
{\small 新编题：活动结束后捡到背包，写招领启事，说明发现地点、领取方式和认领核实要求。}\par
{\small\textbf{行文路线：}物品与地点 → 保管、核实及领取要求 → 帮助转告}\par
\par\vspace{4pt}{\bfseries 可替换模板}\par\nopagebreak
\begin{writing}
{\centering FOUND\\{\bfseries 【拾得物品】}\par}
{\bfseries 【拾得物品，含冠词】} was found {\bfseries 【拾得地点】} after {\bfseries 【拾得时间或活动】}. It has been handed to {\bfseries 【保管与领取地点】} for safekeeping.

{\bfseries 【概括物品内容，完整句】} To make sure it is returned to the right person, the owner should {\bfseries 【核实方式，动词原形】}. Please bring {\bfseries 【领取凭证】} when collecting it. The desk is open {\bfseries 【领取时间】}.

If you know someone who is missing {\bfseries 【相应物品，含冠词】}, please pass on this message. Thank you for your help.

\noindent {\bfseries 【题目要求的署名或发布单位】}\par
\end{writing}
\par\vspace{4pt}{\bfseries 完整套用示例}\par\nopagebreak
{\small 正文 106 词；标题／称呼／落款等另 9 词}\par
\begin{example}
{\centering FOUND\\Backpack at the Student Centre\par}
A backpack was found beside the entrance to the student centre after yesterday's music workshop. It has been handed to the service desk for safekeeping.

The bag contains study materials and a few personal items. To make sure it is returned to the right person, the owner should describe its colour and identify at least two things inside. Please bring a student card or another form of identification when collecting it. The desk is open from nine in the morning to five in the afternoon on weekdays.

If you know someone who is missing a backpack, please pass on this message. Thank you for your help.

\noindent Student Service Centre\par
\end{example}
\par\vspace{4pt}{\bfseries 好词好句：优先记这三处}\par\nopagebreak
{\small\textbf{for safekeeping} — 为妥善保管}\par
{\small\textbf{a form of identification} — 一种身份证明}\par
{\small\textbf{pass on this message} — 转告此消息}\par
{\small\textbf{核对重点：}标题用 FOUND；不把所有可核实细节都公开，也不要把招领写成自己在找失物。}\par
{\fontsize{8}{10}\selectfont\textbf{来源线索：}S05、S18、S19；告示功能迁移重写。}\par
\clearpage\section*{A24  倡议书}\addcontentsline{toc}{section}{A24  倡议书}
\markboth{A24 倡议}{}
{\small 新编题：学生会倡议同学减少食堂食物浪费，说明问题并提出可执行的行动。}\par
{\small\textbf{行文路线：}问题 → 个人与集体可行行动 → 呼吁参与}\par
\par\vspace{4pt}{\bfseries 可替换模板}\par\nopagebreak
\begin{writing}
{\centering {\bfseries 【与倡议目标相符的标题】}\par}
{\bfseries 【问题现象，完整句】} We invite everyone to help improve the situation through a few manageable changes.

Start by {\bfseries 【行动一，动名词】} whenever you {\bfseries 【具体场景，动词短语】}. This can reduce {\bfseries 【直接问题】} without {\bfseries 【不必要的负担，动名词】}. You can also {\bfseries 【行动二，动词原形】} and encourage friends to do the same. If you notice {\bfseries 【具体障碍】}, share your suggestions with {\bfseries 【合适的组织】} so that practical improvements can be made. Even a small change is useful when it becomes a regular habit.

Let us begin with {\bfseries 【明确起步行动】}. By making careful choices in everyday situations, we can help {\bfseries 【共同目标，动词原形】}.

\noindent {\bfseries 【题目要求的署名或发布单位】}\par
\end{writing}
\par\vspace{4pt}{\bfseries 完整套用示例}\par\nopagebreak
{\small 正文 109 词；标题／称呼／落款等另 8 词}\par
\begin{example}
{\centering A Call to Reduce Food Waste\par}
Too much food is left on trays after meals in our dining hall. We invite everyone to help reduce this waste through a few manageable changes.

Choose a smaller portion when you are unsure how much you can finish, and return for more if necessary. When eating with friends, discuss quantities before ordering shared dishes. Students can also tell the dining hall which portion sizes would better suit their needs. These actions save food without requiring anyone to skip a meal or eat less than they need.

Let us begin with our next meal and encourage one another to make careful choices. Every finished portion is a useful step.

\noindent Students' Union\par
\end{example}
\par\vspace{4pt}{\bfseries 好词好句：优先记这三处}\par\nopagebreak
{\small\textbf{manageable changes} — 容易实行的改变}\par
{\small\textbf{portion sizes} — 每份食物的分量}\par
{\small\textbf{make careful choices} — 认真作出选择}\par
{\small\textbf{核对重点：}倡议要让读者知道现在能做什么；区别于报名通知，不能只写时间地点。若题目指定信件格式，另加称呼。}\par
{\fontsize{8}{10}\selectfont\textbf{来源线索：}S18、S19；行动与意义功能综合重写。}\par
\clearpage\section*{A25  会议纪要}\addcontentsline{toc}{section}{A25  会议纪要}
\markboth{A25 纪要}{}
{\small 新编题：记录读书会会议：讨论场地，决定使用图书馆；Chen 负责海报、Wang 负责报名，下周一前完成，周三复核。}\par
{\small\textbf{行文路线：}会议事实 → 决议、负责人和期限 → 后续跟进}\par
\par\vspace{4pt}{\bfseries 可替换模板}\par\nopagebreak
\begin{writing}
{\centering Minutes of {\bfseries 【会议名称】}\par}
\noindent Date: {\bfseries 【日期】}   Place: {\bfseries 【地点】}\\Present: {\bfseries 【与会者】}\par
The committee met to finalise arrangements for {\bfseries 【事项】}. Members discussed {\bfseries 【议题一】}, {\bfseries 【议题二】} and {\bfseries 【议题三】}.

After considering {\bfseries 【备选内容】}, the group agreed to {\bfseries 【决议，动词原形】} because {\bfseries 【理由，完整分句】}. {\bfseries 【负责人一】} will {\bfseries 【任务一，动词原形】}, while {\bfseries 【负责人二】} will {\bfseries 【任务二，动词原形】}. Both tasks must be completed by {\bfseries 【期限】}. Members also agreed that {\bfseries 【另一决议，完整分句】}. {\bfseries 【指定人员】} will collect the necessary information and circulate it to the team before work begins.

The chair will review progress {\bfseries 【检查时间】}. Any difficulties should be reported in advance so that the arrangements can be adjusted.

\noindent Recorded by: {\bfseries 【记录人】}\par
\end{writing}
\par\vspace{4pt}{\bfseries 完整套用示例}\par\nopagebreak
{\small 正文 107 词；标题／称呼／落款等另 23 词}\par
\begin{example}
{\centering Minutes of the Reading Club Meeting\par}
\noindent Date: 12 May 2026\\Place: Room 201\\Present: Li Ming, Chen and Wang\par
The committee met to finalise arrangements for the campus reading event. Members discussed the venue, registration process and publicity materials.

After comparing two available rooms, the group agreed to use the library seminar room because it offered enough seating and a quiet setting. Chen will prepare the poster, while Wang will collect registrations and check participant numbers. Both tasks must be completed by next Monday. Members also agreed that each discussion group should have a student facilitator to keep the conversation focused.

The chair will review progress at a short meeting on Wednesday. Any difficulties should be reported in advance so that the arrangements can be adjusted.

\noindent Recorded by: Li Ming\par
\end{example}
\par\vspace{4pt}{\bfseries 好词好句：优先记这三处}\par\nopagebreak
{\small\textbf{finalise arrangements} — 确定最终安排}\par
{\small\textbf{a student facilitator} — 学生讨论引导员}\par
{\small\textbf{review progress} — 检查进展}\par
{\small\textbf{核对重点：}讨论用过去时；待办可用 will。纪要记录已形成的决议，不写邀请、个人感想或未在会议中决定的内容。}\par
{\fontsize{8}{10}\selectfont\textbf{来源线索：}S08、S16、S18、S19、S29；纪要结构重写。}\par
\clearpage\section*{A26  备忘录}\addcontentsline{toc}{section}{A26  备忘录}
\markboth{A26 备忘录}{}
{\small 新编题：向活动团队说明房间更换，列出需要完成的工作、负责人和确认期限。}\par
{\small\textbf{行文路线：}变更事实 → 操作要求和分工 → 确认时限}\par
\par\vspace{4pt}{\bfseries 可替换模板}\par\nopagebreak
\begin{writing}
\noindent MEMO\\To: {\bfseries 【收件对象】}   From: {\bfseries 【发件人】}\\Date: {\bfseries 【日期】}   Subject: {\bfseries 【主题】}\par
{\bfseries 【活动名称】} has been moved to {\bfseries 【新地点】} because {\bfseries 【变动原因，完整分句】}. The date and starting time remain unchanged.

Please {\bfseries 【需要完成的工作，动词原形】}. {\bfseries 【明确分工，完整句】} We also need to {\bfseries 【通知办法，动词原形】}. Anyone who cannot receive email should be contacted by phone.

Please confirm completion of your tasks by {\bfseries 【确认期限】}. Let me know immediately if anything may delay the preparations, so we can arrange additional help.

\end{writing}
\par\vspace{4pt}{\bfseries 完整套用示例}\par\nopagebreak
{\small 正文 106 词；标题／称呼／落款等另 14 词}\par
\begin{example}
\noindent MEMO\\To: Workshop Team\\From: Li Ming\\Date: 12 May 2026\\Subject: Room Change\par
The reading workshop has been moved to Room 305 because the original room is undergoing repairs. The date and starting time remain unchanged.

Please update the registration page and replace the signs at the building entrance. Chen should check that the new room has enough chairs, while Wang should test the projector before Thursday afternoon. We also need to email registered participants with the new location and a simple map. Anyone who cannot receive email should be contacted by phone.

Please confirm completion of your tasks by four on Thursday. Let me know immediately if anything may delay the preparations, so we can arrange additional help.

\end{example}
\par\vspace{4pt}{\bfseries 好词好句：优先记这三处}\par\nopagebreak
{\small\textbf{remain unchanged} — 保持不变}\par
{\small\textbf{registered participants} — 已报名参加者}\par
{\small\textbf{confirm completion} — 确认已完成}\par
{\small\textbf{核对重点：}To／From／Date／Subject 四项醒目；备忘录通常不加 Dear 和 Yours，与记录过去会议的 minutes 区分。}\par
{\fontsize{8}{10}\selectfont\textbf{来源线索：}S19 第9页、S29；备忘录体裁与安排功能综合新编。}\par
\clearpage\section*{B01  静态构成：饼图／单选比例}\addcontentsline{toc}{section}{B01  静态构成：饼图／单选比例}
\markboth{B01 大作文}{}
{\small 新编数据：某校学生最偏好的阅读方式（单选）：纸质书 45\%，电子书 35\%，有声书 20\%。要求概述数据并评论。}\par
{\small\textbf{行文路线：}对象与构成 → 最大项及关键对比 → 用途差异解释 → 多样化供给}\par
\par\vspace{4pt}{\bfseries 可替换模板}\par\nopagebreak
\begin{writing}
The {\bfseries 【图表类型】} presents how {\bfseries 【调查对象】} are distributed across {\bfseries 【类别】}. {\bfseries 【最大项】} has the largest share, at {\bfseries 【数值】}, followed by {\bfseries 【第二项及数值】}. The figure for {\bfseries 【较小项】} is {\bfseries 【数值】}. These differences show a clear preference for {\bfseries 【主要特征】} within the group surveyed.

One possible reason is that {\bfseries 【原因一，完整分句】}. For people who {\bfseries 【特定需要，动词短语】}, this can make {\bfseries 【选择】} easier or more useful. Another consideration is {\bfseries 【因素二】}, since {\bfseries 【作用机制，完整分句】}. The smaller categories may still meet needs that the leading option cannot satisfy. Their shares should therefore be understood in relation to {\bfseries 【具体情境】}, rather than treated as evidence that they have little value.

{\bfseries 【行动主体】} could respond by {\bfseries 【具体措施，动名词】}. Giving people access to {\bfseries 【相关资源】} would help them make choices that fit their circumstances. The practical aim is to {\bfseries 【贴题目标，动词原形】}.

\end{writing}
\par\vspace{4pt}{\bfseries 完整套用示例}\par\nopagebreak
{\small 正文 150 词}\par
\begin{example}
Reading preferences are divided among three formats in the pie chart. Printed books lead with 45\% of the students surveyed, compared with 35\% for electronic books and 20\% for audiobooks. Print is the largest category, although digital formats together attract more than half of the respondents.

The pattern may reflect the different situations in which people read. A printed book can offer a focused experience away from screen notifications. Electronic books, meanwhile, are easy to carry and search, while audio allows reading to fit around other activities. These advantages help explain why no single format attracts everyone. Preference may depend on the task as much as on the reader.

Libraries could respond by offering a balanced collection and helping students choose suitable formats. The aim should be to make useful reading easier, with different resources supporting different needs. A popular format is valuable when it serves a clear reading purpose.

\end{example}
\par\vspace{4pt}{\bfseries 好词好句：优先记这三处}\par\nopagebreak
{\small\textbf{the students surveyed} — 受调查的学生；不要写全部大学生}\par
{\small\textbf{a focused experience} — 专注的体验}\par
{\small\textbf{a balanced collection} — 均衡配置的馆藏}\par
{\small\textbf{核对重点：}45\% 是最大项，不是 majority；只有互斥且同一总体的类别才能相加。这里电子书与有声书合计55\%。}\par
{\fontsize{8}{10}\selectfont\textbf{来源线索：}S01、S04、S11、S17、S35；静态对比与需求分析综合重写。}\par
\clearpage\section*{B02  静态多选：柱状图／调查动机}\addcontentsline{toc}{section}{B02  静态多选：柱状图／调查动机}
\markboth{B02 大作文}{}
{\small 新编数据：居民选择社区课程时的考虑因素，可多选：时间方便 68\%，实用内容 56\%，低费用 42\%，结交朋友 27\%。要求描述并评论。}\par
{\small\textbf{行文路线：}按选择率排序 → 说明多选口径 → 解释主要因素 → 对应措施}\par
\par\vspace{4pt}{\bfseries 可替换模板}\par\nopagebreak
\begin{writing}
The bar chart reports the factors that {\bfseries 【调查对象】} consider when {\bfseries 【行为，动名词】}. {\bfseries 【最高项】} is selected most often, at {\bfseries 【数值】}, followed by {\bfseries 【第二项及数值】}. {\bfseries 【另外的因素及数据，完整句】} Since respondents could give more than one answer, the percentages refer to overlapping choices.

The leading factor may be important because {\bfseries 【原因一，完整分句】}. Even when {\bfseries 【某种优势存在，完整分句】}, people may hesitate if {\bfseries 【具体阻碍，完整分句】}. The second factor adds a different consideration: {\bfseries 【原因二，完整分句】}. Together, these concerns suggest that people assess {\bfseries 【决策对象】} through several practical needs. A lower selection rate does not make a factor irrelevant to every respondent.

To respond to these needs, {\bfseries 【主体】} could {\bfseries 【措施一，动词原形】} and {\bfseries 【措施二，动词原形】}. Such changes would focus resources on the concerns that appear most often in the survey, while leaving room for individual differences.

\end{writing}
\par\vspace{4pt}{\bfseries 完整套用示例}\par\nopagebreak
{\small 正文 147 词}\par
\begin{example}
Convenient timing is the most frequently selected reason for choosing community courses, mentioned by 68\% of respondents. Practical content follows at 56\%, while low fees and opportunities to meet people receive 42\% and 27\%. Since respondents could select several reasons, the percentages describe overlapping preferences rather than parts of a single total.

These results suggest that a course must first fit into people's daily lives. An interesting class may still be difficult to attend if it conflicts with work or family duties. Practical content can provide an additional reason to participate because learners can use what they learn outside the classroom. Price also matters, though it is selected less often than timing or content.

Organisers could therefore offer several time slots and explain the practical value of each course. Keeping fees reasonable would further improve access. Such arrangements would address the concerns actually reported in the survey.

\end{example}
\par\vspace{4pt}{\bfseries 好词好句：优先记这三处}\par\nopagebreak
{\small\textbf{overlapping preferences} — 可重叠的偏好}\par
{\small\textbf{conflict with work or family duties} — 与工作或家庭责任冲突}\par
{\small\textbf{offer several time slots} — 提供多个时段}\par
{\small\textbf{核对重点：}多选合计193\%并非数据错；不能写剩余比例，也不能把68\%与56\%相加后称为124\%的人。}\par
{\fontsize{8}{10}\selectfont\textbf{来源线索：}S09 第5页、S17、S35；静态多选口径综合新编。}\par
\clearpage\section*{B03  静态表格：不同群体的对比}\addcontentsline{toc}{section}{B03  静态表格：不同群体的对比}
\markboth{B03 大作文}{}
{\small 新编数据：三个年龄组对某社区服务满意／一般／不满意的比例：18—29岁为65／25／10\%，30—49岁为50／30／20\%，50岁及以上为70／20／10\%。要求描述并评论。}\par
{\small\textbf{行文路线：}横向比较主要指标 → 纵向补充另一指标 → 群体需求差异 → 精准改善}\par
\par\vspace{4pt}{\bfseries 可替换模板}\par\nopagebreak
\begin{writing}
The table compares {\bfseries 【指标】} among {\bfseries 【群体】}. {\bfseries 【最高组】} records the highest figure, at {\bfseries 【数值】}, while {\bfseries 【最低组】} has the lowest, at {\bfseries 【数值】}. A further contrast concerns {\bfseries 【另一指标】}: {\bfseries 【关键数据对比，完整句】}. This is a comparison between groups, rather than a record of change over time.

Different circumstances may explain part of the gap. People in {\bfseries 【群体一】} may need {\bfseries 【具体需要】}, whereas those in {\bfseries 【群体二】} may place more value on {\bfseries 【另一需要】}. For example, {\bfseries 【与群体差异相关的情境，完整分句】}. These considerations could affect how people assess {\bfseries 【调查事项】}. Further feedback would be needed to determine which explanation matters most; group averages do not describe every individual.

A useful response would be to {\bfseries 【具体措施，动词原形】}. By asking what each group needs, {\bfseries 【主体】} could improve {\bfseries 【相关服务】} while preserving features that users already value.

\end{writing}
\par\vspace{4pt}{\bfseries 完整套用示例}\par\nopagebreak
{\small 正文 150 词}\par
\begin{example}
Satisfaction with the community service varies across the three age groups in the table. The figure is highest among residents aged fifty and above, at 70\%, and lowest among those aged thirty to forty-nine, at 50\%. The youngest group records 65\%. Dissatisfaction reaches 20\% in the middle group, twice the level in either of the other groups.

Different needs may help explain this contrast. Residents balancing employment and family care may find standard service hours less convenient. Older users may value regular personal assistance, while younger adults may prefer quick online access. These are possible explanations, and further feedback would be needed to identify the actual problems. The figures alone do not reveal individual experiences.

A useful response would be to ask each group what needs improvement and test flexible service hours. Improvements should address specific difficulties while preserving features that users already value, rather than treating all residents alike.

\end{example}
\par\vspace{4pt}{\bfseries 好词好句：优先记这三处}\par\nopagebreak
{\small\textbf{twice the level in ...} — 是……水平的两倍}\par
{\small\textbf{preserving features that users already value} — 保留使用者已经认可的特点}\par
{\small\textbf{treating all residents alike} — 以同一方式对待所有居民}\par
{\small\textbf{核对重点：}年龄组之间的比较不代表同一批人随年龄增长的变化；不同分母的百分比不可直接相加为整体满意率。}\par
{\fontsize{8}{10}\selectfont\textbf{来源线索：}S01、S09、S17、S35；表格与群体对比综合重写。}\par
\clearpage\section*{B04  动态单指标：上升／下降}\addcontentsline{toc}{section}{B04  动态单指标：上升／下降}
\markboth{B04 大作文}{}
{\small 新编数据：某市公共交通出行占比，2020年30\%，2022年36\%，2024年45\%。要求概述趋势并评论。}\par
{\small\textbf{行文路线：}时间、指标和增量 → 两个解释机制 → 持续改善条件}\par
\par\vspace{4pt}{\bfseries 可替换模板}\par\nopagebreak
\begin{writing}
The {\bfseries 【图表类型】} records a clear change in {\bfseries 【指标】} from {\bfseries 【起年】} to {\bfseries 【止年】}. The figure {\bfseries 【rose/fell】} from {\bfseries 【起点数据】} to {\bfseries 【终点数据】}, with {\bfseries 【关键中点或阶段，名词短语】}. Overall, the data indicate {\bfseries 【主要趋势，名词短语】} during the period shown.

One possible explanation is that {\bfseries 【原因一，完整分句】}. This could make it easier for {\bfseries 【群体】} to {\bfseries 【行为，动词原形】}, particularly when {\bfseries 【具体情境，完整分句】}. Another relevant factor may be {\bfseries 【因素二】}. By {\bfseries 【作用方式，动名词】}, it can influence the choices people make in everyday life. These mechanisms offer a plausible interpretation, although the chart itself does not establish which factor caused the change.

The practical response should be to {\bfseries 【对应措施，动词原形】}. If {\bfseries 【必要条件，完整分句】}, {\bfseries 【合理预期，完整分句】}. Attention to these conditions would be more useful than assuming that the same pattern must continue indefinitely.

\end{writing}
\par\vspace{4pt}{\bfseries 完整套用示例}\par\nopagebreak
{\small 正文 150 词}\par
\begin{example}
Public transport accounted for a growing share of journeys in the city between 2020 and 2024. The proportion increased from 30\% to 36\% in 2022 and reached 45\% by 2024, a total gain of fifteen percentage points. The increase was larger in the second half of the period.

Several practical factors may have encouraged this shift. More convenient routes could make public transport a realistic choice for journeys that once required a car. Reliable arrival information may also reduce uncertainty and make transfers easier to plan. Lower travel costs would provide another incentive, especially for people who commute regularly. These possibilities link the trend to everyday decisions, although the chart itself does not establish the causes.

To sustain the improvement, the city should focus on reliability, connections and access for people with limited mobility. A service that works well in ordinary situations is more likely to remain useful over time.

\end{example}
\par\vspace{4pt}{\bfseries 好词好句：优先记这三处}\par\nopagebreak
{\small\textbf{a total gain of ... percentage points} — 共增加……个百分点}\par
{\small\textbf{provide another incentive} — 提供另一种动力}\par
{\small\textbf{people with limited mobility} — 行动不便者}\par
{\small\textbf{核对重点：}30\%到45\%是上升15个百分点，也相当于相对增长50\%；下降题须同步改首句 growing、第二段 shift 与结尾 improvement 的评价。}\par
{\fontsize{8}{10}\selectfont\textbf{来源线索：}S01、S04、S17、S26、S30；动态描述与机制分析综合重写。}\par
\clearpage\section*{B05  动态多指标：升降对照／结构变化}\addcontentsline{toc}{section}{B05  动态多指标：升降对照／结构变化}
\markboth{B05 大作文}{}
{\small 新编数据：某校毕业去向，2020年直接就业60\%、继续升学25\%、其他15\%；2025年分别为45\%、40\%、15\%。要求描述并评论。}\par
{\small\textbf{行文路线：}分别写升降和不变 → 概括差距缩小 → 解释选择差异 → 给出分情境建议}\par
\par\vspace{4pt}{\bfseries 可替换模板}\par\nopagebreak
\begin{writing}
The chart compares {\bfseries 【多个指标】} in {\bfseries 【起年】} and {\bfseries 【止年】}. The share of {\bfseries 【A】} fell from {\bfseries 【数值】} to {\bfseries 【数值】}, while that of {\bfseries 【B】} rose from {\bfseries 【数值】} to {\bfseries 【数值】}. {\bfseries 【C】} remained at {\bfseries 【数值】}. The gap between {\bfseries 【A】} and {\bfseries 【B】} therefore {\bfseries 【narrowed/widened】} over the period.

The shift may reflect changes in {\bfseries 【具体需求】}. {\bfseries 【B】} can offer {\bfseries 【优势】}, which matters to people who {\bfseries 【相关情境，动词短语】}. At the same time, {\bfseries 【A】} remains useful because {\bfseries 【原因，完整分句】}. Different costs and benefits can lead people towards different options, so a larger share does not necessarily make one choice best for everyone. The figures describe the distribution of choices, rather than the reasons for each decision.

{\bfseries 【主体】} should help people compare {\bfseries 【相关选项】} using clear information about {\bfseries 【比较标准】}. Advice based on individual circumstances would make these choices more informed and realistic.

\end{writing}
\par\vspace{4pt}{\bfseries 完整套用示例}\par\nopagebreak
{\small 正文 150 词}\par
\begin{example}
The chart shows a change in graduates' plans between 2020 and 2025. Direct employment fell from 60\% to 45\%, while further study rose from 25\% to 40\%. Other choices remained at 15\%. Employment was still the largest category in 2025, but its lead over further study narrowed from thirty-five to five percentage points.

One possible explanation is that more graduates see additional education as useful preparation for specialised work. They may also want time to explore their interests before committing to a career. Meanwhile, direct employment continues to offer immediate income and practical experience. The changing shares therefore need not mean that one route is suitable for everyone. Costs, interests and career requirements can lead to different reasonable choices.

Universities should provide information on both employment and further study, including their demands and likely benefits. Advice based on individual circumstances would help students choose a route they can realistically pursue.

\end{example}
\par\vspace{4pt}{\bfseries 好词好句：优先记这三处}\par\nopagebreak
{\small\textbf{its lead ... narrowed} — 其领先幅度缩小}\par
{\small\textbf{commit to a career} — 确定职业方向}\par
{\small\textbf{individual circumstances} — 个人具体情况}\par
{\small\textbf{核对重点：}占比下降不能推出就业人数下降；没有毕业生总人数，就只讨论份额。最大类别与增长最快类别也不是一回事。}\par
{\fontsize{8}{10}\selectfont\textbf{来源线索：}S01、S04、S17、S26、S30、S35；多指标对比综合重写。}\par
\clearpage\section*{B06  波动趋势：先升后降／转折}\addcontentsline{toc}{section}{B06  波动趋势：先升后降／转折}
\markboth{B06 大作文}{}
{\small 新编数据：某公共图书馆年度到访人次，2021年12万、2022年18万、2023年16万、2024年20万。要求概述并评论。}\par
{\small\textbf{行文路线：}按时间写关键转折 → 概括总趋势 → 谨慎解释 → 用反馈改善}\par
\par\vspace{4pt}{\bfseries 可替换模板}\par\nopagebreak
\begin{writing}
The figure for {\bfseries 【指标】} changed unevenly between {\bfseries 【起年】} and {\bfseries 【止年】}. It first {\bfseries 【rose/fell】} from {\bfseries 【数据】} to {\bfseries 【数据】}, then {\bfseries 【转折动词短语】}, before reaching {\bfseries 【最终值】}. Despite the temporary {\bfseries 【rise/decline】}, the final level was {\bfseries 【higher/lower】} than the starting point.

Both longer trends and short-term conditions may have shaped this pattern. {\bfseries 【具体因素】} could affect whether people {\bfseries 【行为，动词短语】} in a particular year. Meanwhile, {\bfseries 【另一因素】} may influence demand over a longer period. A single turning point should therefore be interpreted alongside the whole series. It does not by itself show that {\bfseries 【过度推断的判断，完整分句】}. Looking at the context would help distinguish a temporary interruption from a more lasting change.

{\bfseries 【主体】} could compare the figures with {\bfseries 【相关记录】} and gather feedback before adjusting {\bfseries 【相关安排】}. The priority should be to understand the variation and improve {\bfseries 【贴题目标】}.

\end{writing}
\par\vspace{4pt}{\bfseries 完整套用示例}\par\nopagebreak
{\small 正文 155 词}\par
\begin{example}
Annual visits to the public library followed an uneven upward path from 2021 to 2024. They rose from 120,000 to 180,000, slipped to 160,000 in 2023, and then reached 200,000 in the final year. Despite the temporary decline, the final figure was well above the starting level and was the highest recorded.

This pattern may reflect both growing demand and changes in how easily people could use the library. Attractive events can bring in more visitors, while opening hours and access conditions may affect attendance from year to year. A temporary fall should therefore be interpreted alongside the longer pattern. It does not by itself show that residents have lost interest in reading.

The library could compare attendance with its programme schedule and collect visitors' feedback before making changes. Maintaining dependable services and offering useful activities would be a sensible response. The priority is to understand the variation and improve the experience behind the numbers.

\end{example}
\par\vspace{4pt}{\bfseries 好词好句：优先记这三处}\par\nopagebreak
{\small\textbf{an uneven upward path} — 有波动的上升过程}\par
{\small\textbf{slipped to} — 小幅回落至}\par
{\small\textbf{the highest recorded} — 图中记录的最高值}\par
{\small\textbf{核对重点：}有中途下降就不能写 rose steadily。此处统计到访人次，不是独立访客人数，不能把 visits 换成 visitors。}\par
{\fontsize{8}{10}\selectfont\textbf{来源线索：}S01、S17、S26；转折与统计对象辨别综合新编。}\par
\clearpage\section*{B07  双图表：规模与使用情况}\addcontentsline{toc}{section}{B07  双图表：规模与使用情况}
\markboth{B07 大作文}{}
{\small 新编数据：图一，某市公园数量2019年40个、2024年60个；图二，2024年居民每周去公园次数，至少3次50\%、1—2次35\%、不足1次15\%。要求概述两图并评论。}\par
{\small\textbf{行文路线：}每幅图各取核心信息 → 并列联系 → 可能机制与条件 → 具体措施}\par
\par\vspace{4pt}{\bfseries 可替换模板}\par\nopagebreak
\begin{writing}
The two charts examine {\bfseries 【共同主题】} from different angles. The first shows that {\bfseries 【图一关键数据，完整分句】}. In the second, {\bfseries 【图二关键数据，完整分句】}. Together, they present {\bfseries 【并列概括，名词短语】} within the periods and groups shown.

A possible link is that {\bfseries 【合理机制，完整分句】}. When {\bfseries 【具体条件，完整分句】}, people may find it easier to {\bfseries 【行为，动词原形】}. {\bfseries 【补充因素】} can also shape their experience, so the number of {\bfseries 【相关事物】} is only part of the picture. The charts do not by themselves establish that one pattern caused the other. Differences in {\bfseries 【统计范围或相关条件】} should be considered before drawing a stronger conclusion.

{\bfseries 【行动主体】} could use these findings to {\bfseries 【具体行动，动词原形】}. Attention to both {\bfseries 【供给或规模方面】} and {\bfseries 【使用或质量方面】} would help turn the observed changes into practical benefits for the people concerned.

\end{writing}
\par\vspace{4pt}{\bfseries 完整套用示例}\par\nopagebreak
{\small 正文 150 词}\par
\begin{example}
The two charts describe the provision and use of public parks in a city. Park numbers increased from forty in 2019 to sixty in 2024. In the survey for 2024, half of residents reported visiting a park at least three times a week, while only 15\% went less than once a week.

Together, the charts show expanded provision alongside frequent use. More nearby parks may make outdoor activity easier to fit into everyday routines. Safe paths, comfortable seating and space for exercise can also make visits appealing to different age groups. However, the data do not establish that the increase in park numbers caused the reported visiting pattern. Location and service quality may matter as much as the total number.

Future planning should therefore consider where parks are needed and how existing spaces are used. Listening to residents would help turn additional facilities into accessible places that support regular activity.

\end{example}
\par\vspace{4pt}{\bfseries 好词好句：优先记这三处}\par\nopagebreak
{\small\textbf{provision and use} — 设施供给与使用情况}\par
{\small\textbf{alongside frequent use} — 与频繁使用并存}\par
{\small\textbf{as much as the total number} — 与总数量同样重要}\par
{\small\textbf{核对重点：}第二幅图只有2024年，不能写居民访问频率“上升”。两图相关不能直接证明因果；先读清各图年份与单位。}\par
{\fontsize{8}{10}\selectfont\textbf{来源线索：}S01、S15、S17；双图综合与谨慎联系重写。}\par
\clearpage\section*{B08  图画材料：行为对比与寓意}\addcontentsline{toc}{section}{B08  图画材料：行为对比与寓意}
\markboth{B08 大作文}{}
{\small 新编图画题：两名学生准备同一份报告。一人只复制屏幕文字，另一人边阅读边记问题，并与同伴讨论。要求简述画面、解释寓意并评论。}\par
{\small\textbf{行文路线：}画面事实 → 行为差异与主题 → 机制论证 → 行动}\par
\par\vspace{4pt}{\bfseries 可替换模板}\par\nopagebreak
\begin{writing}
The picture shows {\bfseries 【人物】} {\bfseries 【动作，动名词】}. Nearby, {\bfseries 【另一人物】} {\bfseries 【另一动作，动词短语】}. The contrast between their actions highlights {\bfseries 【主题】} and invites us to consider what produces a useful result.

{\bfseries 【正向行为，动名词】} matters because it enables people to {\bfseries 【直接作用，动词原形】}. For example, when {\bfseries 【具体情境，完整分句】}, a person who {\bfseries 【积极做法，动词短语】} can {\bfseries 【可解释的结果，动词原形】}. By contrast, relying only on {\bfseries 【反面做法，动名词】} may leave {\bfseries 【未解决的问题】}. The difference is not simply one of effort: it also concerns how that effort is directed. A sensible approach connects {\bfseries 【行为】} with {\bfseries 【目标】} and makes improvement possible through practice.

We can apply this idea by {\bfseries 【具体行动，动名词】}. Starting with {\bfseries 【小步骤】} and checking the result would turn an abstract lesson into a habit that is useful in everyday study and work.

\end{writing}
\par\vspace{4pt}{\bfseries 完整套用示例}\par\nopagebreak
{\small 正文 152 词}\par
\begin{example}
Two students are preparing the same report in different ways. One simply copies words from a screen, while the other writes down questions and discusses them with a classmate. Their contrasting actions draw attention to the difference between collecting information and understanding it.

Easy access to material can make learning faster, but speed alone does not ensure understanding. A student who asks why an explanation works is more likely to notice gaps and connect ideas. Discussion adds another useful step: explaining a point to someone else can reveal whether it is actually clear. Copying may produce a finished page, yet it offers little help when the learner must solve an unfamiliar problem.

A practical response is to turn reading into an active process. Students can summarise ideas in their own words, test them with examples and ask questions about unclear parts. These habits give the information a purpose beyond completing an assignment.

\end{example}
\par\vspace{4pt}{\bfseries 好词好句：优先记这三处}\par\nopagebreak
{\small\textbf{notice gaps and connect ideas} — 发现缺口并联系知识}\par
{\small\textbf{in their own words} — 用自己的话}\par
{\small\textbf{beyond completing an assignment} — 不止用于完成作业}\par
{\small\textbf{核对重点：}只在题目实际给图画时使用；不用英语一20分、160—200词口径。画面未出现的心理和身份不要写成事实。}\par
{\fontsize{8}{10}\selectfont\textbf{来源线索：}S01、S03、S11、S15、S21、S28；图画写法按英语二篇幅重写。}\par
\clearpage\section*{B09  文字材料：评述观点}\addcontentsline{toc}{section}{B09  文字材料：评述观点}
\markboth{B09 大作文}{}
{\small 新编材料观点：“Online learning should replace classroom teaching because it is more convenient.” 要求解释你是否同意，并用理由和例子支撑。}\par
{\small\textbf{行文路线：}概括原观点与立场 → 承认合理点 → 理由与例子 → 回扣判断标准}\par
\par\vspace{4pt}{\bfseries 可替换模板}\par\nopagebreak
\begin{writing}
The passage presents the view that {\bfseries 【作者核心主张，完整分句】}. I accept its emphasis on {\bfseries 【合理点】}, but I would limit the argument where {\bfseries 【限制条件，完整分句】}.

The proposal has a clear benefit: it can {\bfseries 【好处，动词原形】}, especially for people who {\bfseries 【具体需要，动词短语】}. However, this benefit does not settle the whole question. {\bfseries 【另一考虑因素】} also matters because {\bfseries 【原因，完整分句】}. For example, {\bfseries 【贴题例子，完整分句】}. In such a situation, {\bfseries 【自己的判断，完整分句】}. The example suggests that we should assess the proposal against its purpose and the circumstances in which it is used.

A reasonable approach would be to {\bfseries 【结合立场的做法，动词原形】} while {\bfseries 【兼顾的事项，动名词】}. This preserves the useful part of the author's idea without extending it beyond what the reasoning supports. The final choice should serve {\bfseries 【核心目标】}.

\end{writing}
\par\vspace{4pt}{\bfseries 完整套用示例}\par\nopagebreak
{\small 正文 153 词}\par
\begin{example}
The passage argues that online learning should replace classroom teaching because it offers greater convenience. I agree that flexibility is valuable, but I do not think convenience alone justifies replacing one form of education with another.

Online lessons allow students to revisit explanations and study when travel is difficult. These benefits can widen access to learning. However, a classroom also provides immediate discussion, shared practice and opportunities for teachers to notice confusion. For example, a student conducting a simple experiment may need direct guidance that a recorded video cannot provide. Different subjects and learners therefore call for different arrangements. The quality of teaching matters alongside the method of delivery.

A better approach is to combine the strengths of both forms where appropriate. Schools can use online resources for preparation and review while keeping classroom time for interaction and practical work. The choice should follow learning needs rather than a general preference for convenience.

\end{example}
\par\vspace{4pt}{\bfseries 好词好句：优先记这三处}\par\nopagebreak
{\small\textbf{justify replacing ...} — 构成替换……的充分理由}\par
{\small\textbf{where appropriate} — 在适宜的情况下}\par
{\small\textbf{the method of delivery} — 教学内容的传递方式}\par
{\small\textbf{核对重点：}必须回应材料真正的主张；如果题目问同意程度，要明说。不要把所有文字题都强套“问题—原因—措施”。}\par
{\fontsize{8}{10}\selectfont\textbf{来源线索：}S01、S11、S15、S17、S21、S28；材料观点与有限同意论证重写。}\par
\clearpage\section*{B10  图画＋图表：复合材料}\addcontentsline{toc}{section}{B10  图画＋图表：复合材料}
\markboth{B10 大作文}{}
{\small 新编复合材料：图画中志愿者帮助老人使用医院自助机；图表显示某社区60岁以上居民预约方式，手机应用55\%、电话30\%、现场15\%。要求描述材料并评论。}\par
{\small\textbf{行文路线：}画面与数据分别概括 → 共同主题 → 价值和障碍 → 保留多渠道}\par
\par\vspace{4pt}{\bfseries 可替换模板}\par\nopagebreak
\begin{writing}
The drawing depicts {\bfseries 【人物】} {\bfseries 【动作，动名词】}, while the chart provides figures on {\bfseries 【相关主题】}. {\bfseries 【核心数据一，完整句】} {\bfseries 【核心数据二，完整句】} Read together, the materials bring attention to {\bfseries 【共同问题】}.

The scene illustrates how {\bfseries 【具体机制，完整分句】}. This can benefit {\bfseries 【相关人群】} by {\bfseries 【作用方式，动名词】}. The figures, however, also suggest that {\bfseries 【另一项合理观察，完整分句】}. People may face different needs or practical barriers, including {\bfseries 【具体障碍】}. The materials help identify the issue, but they do not reveal every person's motivation. Any explanation should therefore remain consistent with what is actually shown.

{\bfseries 【行动主体】} should combine {\bfseries 【措施一，名词或动名词】} with {\bfseries 【措施二，名词或动名词】}. This would make {\bfseries 【目标】} easier to achieve while preserving {\bfseries 【需要兼顾的价值】}. Good arrangements should work for people in different circumstances, including those who require additional support.

\end{writing}
\par\vspace{4pt}{\bfseries 完整套用示例}\par\nopagebreak
{\small 正文 152 词}\par
\begin{example}
In the drawing, a volunteer helps an older patient use a hospital self-service machine. The accompanying chart shows how residents aged over sixty book appointments: 55\% use an app, 30\% telephone the service and 15\% visit in person. Together, the materials highlight both the reach of digital services and the continuing need for support.

An app can save travel and waiting time, especially for people who use smartphones confidently. Yet some residents may find unfamiliar instructions or small text difficult to follow. The volunteer's assistance illustrates one way to bridge that gap. The booking figures show that a substantial minority still choose other channels, although they do not tell us why each person makes that choice.

Hospitals should improve digital interfaces while keeping telephone and personal assistance available. Clear guidance and patient support would allow more people to benefit from convenient services without making a particular technology a condition for receiving care.

\end{example}
\par\vspace{4pt}{\bfseries 好词好句：优先记这三处}\par\nopagebreak
{\small\textbf{the accompanying chart} — 随附图表}\par
{\small\textbf{bridge that gap} — 弥合这一差距}\par
{\small\textbf{a substantial minority} — 占比不小的少数群体}\par
{\small\textbf{核对重点：}app占55\%可称 majority；选择电话或现场合计45\%，不能由此断定45\%的人不会用手机。画面与统计样本范围要分清。}\par
{\fontsize{8}{10}\selectfont\textbf{来源线索：}S01、S15、S17；图画与图表联系综合新编。}\par
\clearpage\section*{表达速查：把好句用在合适的地方}\addcontentsline{toc}{section}{表达速查：把好句用在合适的地方}
\markboth{表达速查}{}
本册每篇的三条表达共108条，既包含实用词组，也包含可迁移句法。下面这些结构优先掌握；每次选与任务有关的两三条，不要把整页连续塞进一篇作文。

\begin{longtable}{L{25mm}L{72mm}L{56mm}}\toprule
\textbf{用途} & \textbf{可迁移表达} & \textbf{如何用}\\\midrule\endhead
建议＋结果 & Start by doing ... This will help you ... & 先给具体操作，再解释效果。\\
需求匹配 & a good match for your interest in ... & 推荐理由要对应对方兴趣。\\
邀请角色 & have you as our guest speaker & 说清对方以什么身份参加。\\
尊重时间 & whether this would fit your schedule & 向教授求助或协调安排。\\
承担责任 & I owe you an apology for doing ... & 明确为哪件事道歉。\\
补救办法 & Until then, I can ... & 给出最终解决前的临时安排。\\
事实核查 & check ... against the original records & 核查姓名、时间和数据。\\
申请价值 & support the team with ... & 把技能与岗位需要相连。\\
数据比较 & its lead over ... narrowed & 写两个指标差距缩小。\\
波动 & an uneven upward path & 总体上升但中间有回落。\\
群体范围 & among the students surveyed & 限制在题目调查样本内。\\
条件性判断 & where appropriate & 在适当条件下，不无条件套用。\\
给出机制 & By doing ..., ... can ... & 分词逻辑主语必须与主句一致。\\
谨慎解释 & One possible explanation is that ... & 说明是假设解释，不是图表证明。\\
有依据的推论 & The example suggests that ... & 先有例子，再给推论。\\
具体改进 & collect users’ feedback before ... & 先收集反馈，再改变安排。\\
\bottomrule\end{longtable}
\par\vspace{4pt}{\bfseries 从“显得高级”转向“说得准确”}\par\nopagebreak
increase／rise 通常比不自然的 ascend 或 augment 更好用；the proportion of... 用于比例，the number of... 用于数量；information、advice、equipment 通常不可数。不用为了避重复把所有 people 替换成 individuals 或把 letter 替换成 missive。

\clearpage\section*{图表与语言：交卷前的核对表}\addcontentsline{toc}{section}{图表与语言：交卷前的核对表}
\markboth{核对表}{}
\begin{longtable}{L{37mm}L{120mm}}\toprule
\textbf{容易混淆处} & \textbf{准确写法或判断}\\\midrule\endhead
30\% → 45\% & rose by 15 percentage points；相对增长为50\%，不是15\%。\\
数量120 → 180 & rose by 60；相对增长50\%；rose to 180 是升至180。\\
45\%为最大项 & the largest share；未超过50\%，不用 a majority。\\
不同群体的数据 & 18—29岁与50岁以上的对比，不是时间上的增长过程。\\
只有起点终点 & 能说 increased，不足以证明 steadily／持续平稳。\\
多选合计超过100\% & 可合理；各类别重叠，不能计算“剩余人群”。\\
份额下跌 & 不能在未知总体时断言数量下跌。\\
人次与人数 & visits 是人次；visitors 是访客人数，两者不互换。\\
两幅图相关 & can suggest a possible link，不等于证明因果。\\
过去数据与当前解释 & 已结束年份的数据用过去时；一般机制可用现在时。\\
副词与形容词 & rose sharply；a sharp rise；不要 a sharply rise。\\
两个并列数据 & A and B were 20\% and 30\%, respectively；顺序必须一一对应。\\
建议信动词 & suggest doing／suggest that you do；不是 suggest you to do。\\
期待与道歉 & look forward to hearing；apologise to somebody for doing。\\
咨询中的从句 & Could you tell me when it begins?；从句不用 when does it begin。\\
单复数 & the number of students is ...；a number of students are ...。\\
\bottomrule\end{longtable}
\par\vspace{4pt}{\bfseries 60秒终检}\par\nopagebreak
逐一划去任务清单上的要点；核对图表标题、对象、年份、单位和数字；检查每段是否有具体信息；核对称呼和 Li Ming；检查第三人称单数、时态、冠词、单复数、拼写；最后删去占位提示。

同类题复习时可以只背一篇完整示例，但必须另练一次换题：先只看中文任务写出结构，再用新数据或新活动填内容。若某句离开示例就不会改，优先把那一句改成自己能准确写出的短句。

\clearpage\section*{文字重合检查：方法、结果与边界}\addcontentsline{toc}{section}{文字重合检查：方法、结果与边界}
\markboth{文字重合检查}{}
为落实“不与任何一个文件相似程度大于50\%”，本册采用公开、可复核的文字重合指标。它衡量英文原文复用，不衡量意义相似度，也不能代表某个商业系统的查重结论。以下结果仅适用于本次版本及已经识别出的文本。

\par\vspace{4pt}{\bfseries 比较范围与计算方法}\par\nopagebreak
处理用户给出的40个PDF，共1164页（含重复文件及408大纲）。可提取的正文使用PDF文字层；扫描正文或乱码文字层经macOS Vision识别。S07使用工作区已保存的完整OCR；S11与S12字节一致，但仍分别纳入比对。S02是计算机考纲，不参与写作内容筛选，仍列入机器比对。图片中的文字和扫描件仍可能存在漏识别或错识别。

对每篇模板与每篇示例分别比较，不靠整本文档长度稀释结果。英文转为小写，标点、数字不参与匹配，连字符词拆开；不删除常见虚词。模板的中文占位提示、各篇题目说明和格式栏不计入正文检查。固定英文部分单独比对，避免用填空示例的新细节稀释模板本身的重合。

主指标：新稿中凡属于与某一来源相同的连续4个英文词的词位，标记为重合；重叠词位只计算一次。覆盖率＝标记词位数÷新稿英文词元总数。对40份来源逐一计算，取最高值。辅助指标用相同方法检查连续5词。不是把与不同来源的命中合并后冒充单一文件相似度。

\begin{longtable}{L{70mm}L{87mm}}\toprule
\textbf{核查项} & \textbf{结果}\\\midrule\endhead
比较单元 & 36篇模板＋36篇完整示例；各自对照40份文件\\
连续4词覆盖率最高值 & 20.63\%\\
连续5词覆盖率最高值 & 10.20\%\\
本次阈值结果 & 72个比较单元的最高4词覆盖率均低于50\%\\
适用边界 & 不等同于语义相似度；OCR错误可能漏掉部分原文重合。\\
\bottomrule\end{longtable}
同一文体的格式、常用句法和考试任务本就会相近。因此本册没有承诺“任何算法都低于50\%”。检查脚本、来源清单及逐篇CSV保存在源文件目录，能够重新执行。

\clearpage\par\vspace{4pt}{\bfseries 逐篇最高覆盖率（分别与每一份来源比较）}\par\nopagebreak
{\small
\begin{longtable}{L{13mm}L{67mm}L{20mm}L{20mm}L{25mm}}\toprule
\textbf{编号} & \textbf{类型} & \textbf{模板4词} & \textbf{示例4词} & \textbf{对应来源}\\\midrule\endhead
A01 & 邀请信 & 7.55\% & 4.08\% & S34/S34\\
A02 & 推荐信：书籍／活动／地点 & 7.35\% & 0.00\% & S34/S01\\
A03 & 推荐信：推荐他人 & 0.00\% & 0.00\% & S01/S01\\
A04 & 建议信 & 12.31\% & 3.96\% & S07/S19\\
A05 & 计划信：提出方案并征求意见 & 4.88\% & 4.90\% & S36/S36\\
A06 & 介绍信：介绍项目并邀请参与 & 12.31\% & 3.85\% & S34/S06\\
A07 & 答复信：回复咨询 & 0.00\% & 4.85\% & S01/S11\\
A08 & 交流信：表达感受并询问详情 & 6.58\% & 4.95\% & S18/S18\\
A09 & 感谢信 & 17.11\% & 8.26\% & S16/S27\\
A10 & 道歉信 & 5.63\% & 3.85\% & S01/S19\\
A11 & 投诉信 & 6.33\% & 7.69\% & S17/S18\\
A12 & 咨询信 & 6.67\% & 4.67\% & S10/S10\\
A13 & 请求信：请求帮助 & 12.82\% & 9.17\% & S16/S16\\
A14 & 申请信：项目／志愿者 & 0.00\% & 3.85\% & S01/S24\\
A15 & 求职信 & 10.96\% & 7.62\% & S10/S10\\
A16 & 祝贺信 & 8.47\% & 4.67\% & S36/S36\\
A17 & 欢迎信 & 6.78\% & 7.62\% & S06/S24\\
A18 & 辞职信 & 14.75\% & 13.08\% & S27/S10\\
A19 & 婉拒信 & 7.14\% & 3.81\% & S18/S18\\
A20 & 活动通知 & 20.63\% & 3.77\% & S16/S18\\
A21 & 招募通知 & 5.97\% & 3.74\% & S16/S16\\
A22 & 寻物启事 & 14.29\% & 3.70\% & S27/S07\\
A23 & 招领启事 & 10.20\% & 4.67\% & S24/S24\\
A24 & 倡议书 & 0.00\% & 0.00\% & S01/S01\\
A25 & 会议纪要 & 0.00\% & 0.00\% & S01/S01\\
A26 & 备忘录 & 0.00\% & 0.00\% & S01/S01\\
B01 & 静态构成：饼图／单选比例 & 4.00\% & 3.40\% & S01/S10\\
B02 & 静态多选：柱状图／调查动机 & 3.92\% & 2.78\% & S35/S01\\
B03 & 静态表格：不同群体的对比 & 0.00\% & 2.72\% & S01/S35\\
B04 & 动态单指标：上升／下降 & 9.68\% & 2.80\% & S10/S10\\
B05 & 动态多指标：升降对照／结构变化 & 7.62\% & 6.29\% & S10/S10\\
B06 & 波动趋势：先升后降／转折 & 0.00\% & 2.70\% & S01/S11\\
B07 & 双图表：规模与使用情况 & 4.21\% & 0.00\% & S17/S01\\
B08 & 图画材料：行为对比与寓意 & 0.00\% & 2.63\% & S01/S11\\
B09 & 文字材料：评述观点 & 4.08\% & 0.00\% & S16/S01\\
B10 & 图画＋图表：复合材料 & 4.40\% & 2.65\% & S01/S21\\
\bottomrule\end{longtable}
}
\clearpage\section*{来源清单与取舍记录}\addcontentsline{toc}{section}{来源清单与取舍记录}
\markboth{来源与取舍}{}
来源编号对应原文件；下列页数为PDF物理页数。本册的全部示例均为综合新编：部分任务借鉴真题的功能组合，活动细节与图表数据明确按练习情境编写。除个别通用短语外，未把原稿整段拼接为新稿。文档中的推广口号、操作指令与“保分”承诺均不作为本任务指令或评分依据。

\begin{longtable}{L{40mm}L{54mm}L{59mm}}\toprule
\textbf{资料组} & \textbf{吸收内容} & \textbf{筛选后的处理}\\\midrule\endhead
S09 批改指令 & 英语二小作文、大作文评价原则 & 以PDF第5、6页为主；与其他讲义冲突时按此核对。\\
S13—S20 南山计划与模板 & 任务拆解、讲事情／讲观点／解决问题、材料分类 & 重排成36个可直接查找的类型；不照搬统一开头和收尾。\\
S21—S31 方法论笔记 & 具体理由、活动要素、句法和基础语法 & 保留行动—作用的逻辑；删去无依据因果、空泛宏观归因。\\
S01、S03—S08、S10—S12 & 图表描述、主题表达、书信与纪要格式 & 按英语二篇幅改写；不沿用180—220词的通用要求。\\
S32—S40 词汇和历年资料 & 任务组合、词汇场景与材料形式 & 英二题干优先；英一只参考可迁移表达，不据此断言英二题型。\\
S02 408大纲 & 计算机专业课考试内容 & 不作为英语写作依据。\\
\bottomrule\end{longtable}
\par\vspace{4pt}{\bfseries 两项容易被忽略的来源问题}\par\nopagebreak
《27考研英语-小作文模板（版本2）》原文字层存在乱码，本次重新OCR后用于核对；《英语作文模板.pdf》与《英语作文模板 (1).pdf》内容和文件哈希一致，属于重复材料。已有《英语二作文模板总纲》也纳入对照，未直接复制为本册。

原资料里既有“约150词”，也有“至少150词”或适用于英语一的篇幅建议；本册不混用。实际作答最终以当场题干为准。自然、准确的短句可以保留，生僻词和长句不自动加分。

\clearpage\par\vspace{4pt}{\bfseries 原文件索引 1—14}\par\nopagebreak
{\small
\begin{longtable}{L{13mm}L{104mm}L{11mm}L{25mm}}\toprule
\textbf{编号} & \textbf{原文件名} & \textbf{页数} & \textbf{读取方式／备注}\\\midrule\endhead
S01 & 1.27考研英语大作文模板（英一英二通用）.pdf & 88 & 文本层\\
S02 & 27考研408大纲.pdf & 20 & 非英语材料\\
S03 & 27考研英语-大作文图画模板.pdf & 5 & 文本层\\
S04 & 27考研英语-大作文图表模板.pdf & 5 & 文本层\\
S05 & 27考研英语-小作文模板（版本1，推荐）.pdf & 9 & 文本层\\
S06 & 27考研英语-小作文模板（版本2）.pdf & 3 & OCR\\
S07 & mmg作文模板.pdf & 49 & OCR\\
S08 & 帕拉迪宇考研英语作文模板.pdf & 56 & 文本层\\
S09 & 考研英语批改指令.pdf & 7 & 评分依据\\
S10 & 英语二作文模板总纲.pdf & 126 & 文本层\\
S11 & 英语作文模板 (1).pdf & 74 & 文本层\\
S12 & 英语作文模板.pdf & 74 & 与S11重复\\
S13 & 27南山作文模板吃鸡计划4.0（英二大作文学习计划）.pdf & 2 & 文本层\\
S14 & 27南山作文模板吃鸡计划4.0（英二小作文学习计划）.pdf & 1 & 文本层\\
\bottomrule\end{longtable}
}
S01—S12 位于工作区“英语二”目录；S13—S40 位于用户指定的“27南山作文模板吃鸡计划4.0（英二）”资料目录。机器清单保留每份原文件的完整路径与SHA-256，便于核对。

\clearpage\par\vspace{4pt}{\bfseries 原文件索引 15—27}\par\nopagebreak
{\small
\begin{longtable}{L{13mm}L{104mm}L{11mm}L{25mm}}\toprule
\textbf{编号} & \textbf{原文件名} & \textbf{页数} & \textbf{读取方式／备注}\\\midrule\endhead
S15 & 27南山作文吃鸡计划大作文万用模版.pdf & 12 & OCR\\
S16 & 27南山作文吃鸡计划小作文万用模版.pdf & 9 & OCR\\
S17 & 27南山作文模板吃鸡计划（大作文+素材库）.pdf & 41 & OCR\\
S18 & 27南山作文模板吃鸡计划（小作文素材库）.pdf & 12 & 文本层\\
S19 & 27南山作文模板吃鸡计划（小作文）.pdf & 43 & OCR\\
S20 & 【补充至90页】27南山大作文第二段补充句.pdf & 1 & 文本层\\
S21 & 27Monkey作文强化课笔记：大作文图画+图表.pdf & 20 & OCR\\
S22 & 27Monkey作文强化课笔记：导学+小作文.pdf & 19 & OCR\\
S23 & 27Monkey作文强化课笔记：思路结构进阶与金句.pdf & 7 & OCR\\
S24 & 27周思成考研英语写作基础课（英二）.pdf & 46 & OCR\\
S25 & 27周思成考研英语写作零基础.pdf & 27 & OCR\\
S26 & 27柴荣老师英语二大作文笔记.pdf & 28 & OCR\\
S27 & 27柴荣老师英语二小作文笔记.pdf & 23 & OCR\\
\bottomrule\end{longtable}
}
S01—S12 位于工作区“英语二”目录；S13—S40 位于用户指定的“27南山作文模板吃鸡计划4.0（英二）”资料目录。机器清单保留每份原文件的完整路径与SHA-256，便于核对。

\clearpage\par\vspace{4pt}{\bfseries 原文件索引 28—40}\par\nopagebreak
{\small
\begin{longtable}{L{13mm}L{104mm}L{11mm}L{25mm}}\toprule
\textbf{编号} & \textbf{原文件名} & \textbf{页数} & \textbf{读取方式／备注}\\\midrule\endhead
S28 & 27王晶婷考研英语写作模板课：大作文.pdf & 57 & OCR\\
S29 & 27王晶婷考研英语写作模板课：小作文.pdf & 19 & OCR\\
S30 & 27石雷鹏大作文技巧精讲（英二）.pdf & 10 & OCR\\
S31 & 27石雷鹏小作文技巧精讲+实训1-3.pdf & 17 & OCR\\
S32 & 【@牛南山】考研作文15类作文主题词汇.pdf & 15 & 文本层\\
S33 & 【南山整理】历年真题范文英一大作文 .pdf & 39 & OCR\\
S34 & 【南山整理】历年真题范文英一小作文 .pdf & 44 & OCR\\
S35 & 【南山整理】历年真题范文英二大作文.pdf & 48 & OCR\\
S36 & 【南山整理】历年真题范文英二小作文 .pdf & 34 & OCR\\
S37 & 英语一大作文真题合集 .pdf & 27 & OCR\\
S38 & 英语一小作文真题合集 .pdf & 17 & OCR\\
S39 & 英语二大作文真题合集 .pdf & 17 & OCR\\
S40 & 英语二小作文真题合集 .pdf & 13 & OCR\\
\bottomrule\end{longtable}
}
S01—S12 位于工作区“英语二”目录；S13—S40 位于用户指定的“27南山作文模板吃鸡计划4.0（英二）”资料目录。机器清单保留每份原文件的完整路径与SHA-256，便于核对。

\end{document}
